% =====================================================================
%  The Idiot Brain webinars - shared Beamer preamble (the banner look)
% ---------------------------------------------------------------------
%  Usage (in slides/webinarNN/webinarNN.tex):
%      \documentclass[aspectratio=169,11pt]{beamer}
%      % =====================================================================
%  The Idiot Brain webinars - shared Beamer preamble (the banner look)
% ---------------------------------------------------------------------
%  Usage (in slides/webinarNN/webinarNN.tex):
%      \documentclass[aspectratio=169,11pt]{beamer}
%      % =====================================================================
%  The Idiot Brain webinars - shared Beamer preamble (the banner look)
% ---------------------------------------------------------------------
%  Usage (in slides/webinarNN/webinarNN.tex):
%      \documentclass[aspectratio=169,11pt]{beamer}
%      % =====================================================================
%  The Idiot Brain webinars - shared Beamer preamble (the banner look)
% ---------------------------------------------------------------------
%  Usage (in slides/webinarNN/webinarNN.tex):
%      \documentclass[aspectratio=169,11pt]{beamer}
%      \input{../common/preamble.tex}
%  Build from the deck folder:  latexmk -lualatex -outdir=build webinarNN.tex
%  Art (emoji, mascot moods, title banner) comes from common/build_art.py.
%
%  All styling lives here, so decks only hold content. Palette and fonts
%  are the banner's (banner/build_banner.py): plum text on cream, pastel
%  accents, Mardoto for text, Glazz for big titles, Twemoji icons and
%  the pink brain mascot (DECISIONS.md #4-#6).
% =====================================================================

% ---------------------------------------------------------------------
% Fonts (LuaLaTeX + fontspec), shared with the banner
% ---------------------------------------------------------------------
\usepackage{fontspec}
\tracinglostchars=3 % a glyph missing from a font stops the build instead of printing nothing
\defaultfontfeatures{Ligatures=TeX}
\def\idiotfontdir{../../banner/assets/fonts/}
\setsansfont{Mardoto}[Path=\idiotfontdir, Extension=.ttf, UprightFont=*-Regular, BoldFont=*-Bold]
\setmainfont{Mardoto}[Path=\idiotfontdir, Extension=.ttf, UprightFont=*-Regular, BoldFont=*-Bold]
\newfontfamily\mediumfont{Mardoto-Medium}[Path=\idiotfontdir, Extension=.ttf]
\newfontfamily\blackfont{Mardoto-Black}[Path=\idiotfontdir, Extension=.ttf]
\newfontfamily\titlefont{GlazzArm-Regular}[Path=\idiotfontdir, Extension=.ttf] % Armenian letters only
\newfontfamily\handfont{Ink Free} % Windows font with Armenian; speech bubbles only
% Mardoto has no italic: emphasis is medium weight in the accent colour instead.
\renewcommand{\emph}[1]{{\mediumfont\color{plum}#1}}
% No hyphenation: split words are hard to read on a shared screen.
\hyphenpenalty=10000
\exhyphenpenalty=10000

% ---------------------------------------------------------------------
% Packages
% ---------------------------------------------------------------------
\usepackage{graphicx}
\usepackage{booktabs}
\usepackage[most]{tcolorbox}
\usepackage{tikz}
\usetikzlibrary{arrows.meta, positioning, calc, shapes.geometric, shapes.callouts, shapes.misc,
  decorations.pathmorphing, decorations.pathreplacing, backgrounds, fit, patterns.meta}
\usepackage{pgfplots}
\pgfplotsset{compat=1.18}
\usepgfplotslibrary{fillbetween}
\graphicspath{{../common/art/}{images/}}

% ---------------------------------------------------------------------
% Colours: the banner palette. Charts with 3+ series use the Armenian
% flag colours (amRed / amBlue / amOrange).
% ---------------------------------------------------------------------
\definecolor{plum}{HTML}{4A2340}
\definecolor{brainpink}{HTML}{F7A8C4}
\definecolor{fold}{HTML}{E57FA6}
\definecolor{cream}{HTML}{FFF6EC}
\definecolor{cardwhite}{HTML}{FFFDF9}
\definecolor{peach}{HTML}{FFD9C2}
\definecolor{lavender}{HTML}{D9C8F5}
\definecolor{mint}{HTML}{B5E8D5}
\definecolor{butter}{HTML}{FFE8A3}
\definecolor{blush}{HTML}{FF8FA8}
\definecolor{amRed}{HTML}{D90012}
\definecolor{amBlue}{HTML}{0033A0}
\definecolor{amOrange}{HTML}{F2A800}
% darker partners for text on cream
\colorlet{mintdark}{mint!45!black!80!green}
\definecolor{mintdark}{HTML}{2E7D62}
\definecolor{lavenderdark}{HTML}{5B3FA0}
\definecolor{peachdark}{HTML}{B5532A}
\definecolor{pinkdark}{HTML}{C2185B}
% names used by the figure files
\colorlet{ink}{plum}
\colorlet{paper}{cream}
\colorlet{mist}{lavender!40!cream}
\definecolor{night}{HTML}{2B2150}
\colorlet{softgray}{plum!60!cream}
% accent of the current part; \partaccent{<colour>} switches it
\colorlet{accent}{butter}
\newcommand{\partaccent}[1]{\colorlet{accent}{#1}}

% ---------------------------------------------------------------------
% Art: \emoji, \mascot (pages of common/art/*.pdf, see build_art.py)
% ---------------------------------------------------------------------
\makeatletter
\newcommand{\artpage}[3]{\expandafter\def\csname idiot@art@#1@#2\endcsname{#3}}
\input{../common/art/art_index.tex}
\newcommand{\idiot@art}[3]{% kind, name, size options
  \ifcsname idiot@art@#1@#2\endcsname
    \includegraphics[#3, page=\csname idiot@art@#1@#2\endcsname]{#1.pdf}%
  \else
    \PackageError{idiot}{Unknown #1 '#2'; add it to build_art.py and rebuild the art}{}%
  \fi}
% \emoji[<height>]{name}: inline icon, sits on the text baseline like a letter
\newcommand{\emoji}[2][1.15em]{\raisebox{-0.2em}{\idiot@art{emoji}{#2}{height=#1}}}
% \bigemoji[<height>]{name}: no baseline shift, for TikZ nodes and big pictures
\newcommand{\bigemoji}[2][1.5cm]{\idiot@art{emoji}{#2}{height=#1}}
% \mascot[<width>]{mood}: panic, happy, scared, sleepy, sick, confused, smug, hungry, thinking
\newcommand{\mascot}[2][2.6cm]{\idiot@art{mascot}{#2}{width=#1}}
\makeatother
% the older figure files place the mascot with \pic{noggin=<mood>}
\tikzset{pics/noggin/.style={code={\node[inner sep=0, transform shape] {\mascot[2.7cm]{#1}};}}}

% ---------------------------------------------------------------------
% Beamer look
% ---------------------------------------------------------------------
\setbeamertemplate{navigation symbols}{}
\setbeamercolor{normal text}{fg=plum, bg=cream}
\setbeamercolor{structure}{fg=plum}
\setbeamercolor{alerted text}{fg=pinkdark}
\setbeamerfont{alerted text}{series=\bfseries}
\setbeamercolor{frametitle}{fg=plum}
\setbeamersize{text margin left=9mm, text margin right=9mm}
\setbeamertemplate{frame begin}{\raggedright} % justified Armenian leaves big gaps

% cream paper with the banner's faint dot grid
\setbeamertemplate{background}{%
  \begin{tikzpicture}
    \fill[cream] (0,0) rectangle (\paperwidth,\paperheight);
    \fill[pattern={Dots[distance=4.4mm, radius=0.28mm]}, pattern color=plum!9!cream]
      (0,0) rectangle (\paperwidth,\paperheight);
  \end{tikzpicture}}

% topic pill, top right of every frame: set by \topic (emoji + short name)
\def\idiottopicemoji{}
\def\idiottopicname{}
\newcommand{\idiotpill}[3][cardwhite]{% fill, emoji, text
  \tikz[baseline=(p.base)]{\node[fill=#1, draw=plum, line width=0.7pt, rounded corners=2.6mm,
      inner xsep=2.2mm, inner ysep=1.1mm, font=\mediumfont\footnotesize, text=plum] (p)
      {\ifx\relax#2\relax\else\emoji[1.05em]{#2}\hspace{1.2mm}\fi#3};
    \begin{scope}[on background layer]
      \fill[plum, rounded corners=2.6mm] ([xshift=0.7mm, yshift=-0.7mm]p.south west) rectangle ([xshift=0.7mm, yshift=-0.7mm]p.north east);
    \end{scope}}}
% \chip{emoji}{text}: inline sticker label, same as the banner's topic chips
\newcommand{\chip}[3][cardwhite]{\idiotpill[#1]{#2}{#3}}

\setbeamertemplate{frametitle}{%
  \vskip 3.5mm
  \begin{minipage}[t]{0.66\paperwidth}
    \raggedright{\blackfont\fontsize{17.5}{20.5}\selectfont\insertframetitle}\par
    \vskip 1.3mm
    \tikz\fill[accent, rounded corners=0.6mm] (0,0) rectangle (1.5cm,1.1mm);
  \end{minipage}%
  \hfill
  \ifx\idiottopicname\empty\else
    \raisebox{-1mm}{\idiotpill{\idiottopicemoji}{\idiottopicname}}%
  \fi
  \vskip 1mm}

\setbeamertemplate{footline}{%
  \hbox to \paperwidth{\hskip 9mm
    {\mediumfont\fontsize{7}{8}\selectfont\color{plum!55!cream}Ապուշ ուղեղը · \insertshorttitle}\hfill
    {\blackfont\fontsize{8}{9}\selectfont\color{plum!70!cream}\insertframenumber}\hskip 9mm}
  \vskip 3mm}

\setbeamertemplate{itemize item}{\tikz[baseline=-0.65ex]\draw[plum, fill=accent, line width=0.6pt] (0,0) circle (0.95mm);}
\setbeamertemplate{itemize subitem}{\tikz[baseline=-0.65ex]\fill[plum!70] (0,0) circle (0.6mm);}
\setbeamertemplate{enumerate item}{\tikz[baseline=(n.base)]\node[fill=accent, draw=plum, line width=0.6pt,
  circle, inner sep=0.5mm, minimum size=4.2mm, font=\blackfont\scriptsize] (n) {\insertenumlabel};}
\setlength{\leftmargini}{5mm}
\setbeamertemplate{itemize/enumerate body begin}{\setlength{\itemsep}{1.6mm}}
\setbeamertemplate{itemize/enumerate subbody begin}{\setlength{\itemsep}{0.6mm}\small}

% ---------------------------------------------------------------------
% Structure: parts (full-colour openers) and topics (book sections)
% ---------------------------------------------------------------------
% \topic{<emoji>}{<short name>}: starts a book section; the pill shows on every following frame
\newcommand{\topic}[2]{\subsection{#2}\gdef\idiottopicemoji{#1}\gdef\idiottopicname{#2}}
\newcommand{\notopic}{\gdef\idiottopicemoji{}\gdef\idiottopicname{}}

% \partopener{<label>}{<title>}{<subtitle>}{<mood>}{<colour>}
\newcommand{\partopener}[5]{%
  \section{#2}\notopic\partaccent{#5}%
  \begin{frame}[plain]
    \begin{tikzpicture}[remember picture, overlay]
      \fill[#5] (current page.south west) rectangle (current page.north east);
      \fill[pattern={Dots[distance=4.4mm, radius=0.28mm]}, pattern color=plum!12!#5]
        (current page.south west) rectangle (current page.north east);
      \fill[cardwhite, opacity=0.45] ([xshift=-3.3cm, yshift=-0.2cm]current page.east) circle (3.1cm);
      \node[anchor=north west, inner sep=0] at ([xshift=10mm, yshift=-12mm]current page.north west)
        {\idiotpill{}{#1}};
      \node[anchor=west, text=plum, align=left, text width=9.2cm, inner sep=0]
        (t) at ([xshift=10mm, yshift=5mm]current page.west)
        {\titlefont\fontsize{38}{40}\selectfont #2};
      \node[anchor=north west, text=plum, text width=9.2cm, inner sep=0, align=left]
        at ([yshift=-5mm]t.south west) {\mediumfont\large #3};
      \node at ([xshift=-3.3cm, yshift=-0.2cm]current page.east) {\mascot[4.6cm]{#4}};
    \end{tikzpicture}
  \end{frame}}

% \topicopener{<emoji>}{<short name>}{<book's section title>}{<what it is about>}{<mood>}
\newcommand{\topicopener}[5]{%
  \topic{#1}{#2}%
  \begin{frame}[plain]
    \begin{tikzpicture}[remember picture, overlay]
      \fill[accent!55!cream] (current page.south west) rectangle (current page.north east);
      \fill[pattern={Dots[distance=4.4mm, radius=0.28mm]}, pattern color=plum!10!accent!55!cream]
        (current page.south west) rectangle (current page.north east);
      \node[anchor=north west, inner sep=0] at ([xshift=10mm, yshift=-12mm]current page.north west)
        {\idiotpill{#1}{#2}};
      \node[anchor=west, text=plum, text width=9.6cm, align=left, inner sep=0]
        (t) at ([xshift=10mm, yshift=6mm]current page.west)
        {\blackfont\fontsize{24}{28}\selectfont «#3»};
      \node[anchor=north west, text=plum, text width=9.6cm, align=left, inner sep=0]
        at ([yshift=-4mm]t.south west) {\mediumfont\large #4};
      \node at ([xshift=-3.0cm, yshift=-0.4cm]current page.east) {\mascot[4.2cm]{#5}};
      \node[rotate=-12] at ([xshift=-1.4cm, yshift=2.3cm]current page.east) {\bigemoji[1.5cm]{#1}};
    \end{tikzpicture}
  \end{frame}}

% ---------------------------------------------------------------------
% Boxes (sticker style: plum outline + hard plum shadow, like the banner)
% ---------------------------------------------------------------------
\tcbset{
  sticker/.style={enhanced, arc=3mm, boxrule=0.9pt, colframe=plum, colback=cardwhite, halign=left,
    left=3mm, right=3mm, top=2mm, bottom=2mm, fontupper=\normalsize, coltext=plum,
    shadow={1.1mm}{-1.1mm}{0mm}{fill=plum}},
  tagged/.style={sticker, colback=#1!45!cream, colbacktitle=#1, coltitle=plum,
    fonttitle=\blackfont\footnotesize, top=3mm,
    attach boxed title to top left={xshift=4mm, yshift*=-2.4mm},
    boxed title style={arc=2.2mm, boxrule=0.9pt, colframe=plum, left=1.5mm, right=1.5mm, top=0.6mm, bottom=0.6mm}},
}
\newtcolorbox{card}[1][]{sticker, #1}
\newtcolorbox{tryit}[1][]{tagged=mint, title={\emoji[1em]{testtube}\ Փորձիր հիմա}, #1}
\newtcolorbox{poll}[1][]{tagged=butter, title={\emoji[1em]{hand}\ Հարցում}, #1}
\newtcolorbox{chatbox}[1][]{tagged=lavender, title={\emoji[1em]{speech}\ Գրիր չաթում}, #1}
\newtcolorbox{scienceupdate}[1][]{tagged=peach, fontupper=\small, title={\emoji[1em]{microscope}\ Գիտությունն այսօր}, #1}
\newtcolorbox{bookquote}[1][]{sticker, colback=cardwhite, fontupper=\mediumfont, left=4mm,
  borderline west={1.4mm}{0pt}{accent},
  before upper={{\blackfont\footnotesize\emoji[1em]{book}\ Բըրնեթը գրում է}\par\smallskip}, #1}

% Brain.exe-style popup window: \begin{popup}{<window title>} ... \end{popup}
\newcommand{\winbuttons}{%
  \tikz[baseline=-0.6ex, line width=0.6pt]{
    \foreach \i/\f in {0/cardwhite, 1/cardwhite, 2/blush}{
      \filldraw[fill=\f, draw=plum, rounded corners=0.6mm] (\i*4.3mm,-1.6mm) rectangle ++(3.3mm,3.3mm);}
    \draw[plum] (0.8mm,-0.8mm) -- (2.5mm,-0.8mm);
    \draw[plum] (5.1mm,-0.8mm) rectangle (6.8mm,0.9mm);
    \draw[plum] (9.4mm,-0.9mm) -- (11.1mm,0.8mm) (9.4mm,0.8mm) -- (11.1mm,-0.9mm);}}
\newtcolorbox{popup}[2][]{sticker, colback=cardwhite, colbacktitle=brainpink, coltitle=plum,
  fonttitle=\blackfont\small, title={\emoji[1.1em]{warning}\hspace{1.5mm}#2\hfill\winbuttons},
  titlerule=0.9pt, toptitle=0.8mm, bottomtitle=0.8mm, #1}

% Myth or fact: \mythorfact{<statement>}{<verdict>}{<colour>}{<explanation>}; verdict shows on slide 2
\newcommand{\mythorfact}[4]{%
  \begin{center}
  \begin{tikzpicture}
    \node[draw=plum, line width=0.9pt, fill=cardwhite, rounded corners=4mm, inner xsep=8mm, inner ysep=6mm,
          text width=10cm, align=center, font=\blackfont\Large, text=plum] (mf) {«#1»};
    \begin{scope}[on background layer]
      \fill[plum, rounded corners=4mm] ([xshift=1.2mm, yshift=-1.2mm]mf.south west) rectangle ([xshift=1.2mm, yshift=-1.2mm]mf.north east);
    \end{scope}
    \node[anchor=west] at ([xshift=6mm]mf.north west) {\idiotpill[butter]{question}{Առասպել, թե՞ փաստ}};
    \onslide<2->{\node[draw=#3, text=#3, line width=2pt, rounded corners=2mm, rotate=-10, font=\blackfont\LARGE,
      inner sep=4pt, fill=cardwhite] at ([xshift=-8mm, yshift=1mm]mf.south east) {#2};}
  \end{tikzpicture}
  \end{center}
  \onslide<2->{#4}}

% Big question with the confused mascot: \bigquestion{<question>}{<small line>}
\newcommand{\bigquestion}[2]{%
  \vfill
  \begin{center}
  \begin{tikzpicture}
    \node[text width=9.6cm, align=center, font=\blackfont\fontsize{24}{29}\selectfont, text=plum] (q) {#1};
    \node[below=4mm of q, text width=9.6cm, align=center, font=\mediumfont\large, text=softgray] {#2};
    \node[anchor=east] at ([xshift=-2mm]q.west) {\mascot[2.6cm]{confused}};
  \end{tikzpicture}
  \end{center}
  \vfill}

% Small source line at the bottom of a frame: \source{Loftus & Palmer, 1974}
\newcommand{\source}[1]{%
  \begin{tikzpicture}[remember picture, overlay]
    \node[anchor=south east, font=\fontsize{6.5}{7.5}\selectfont, text=plum!60!cream, inner sep=0]
      at ([xshift=-17mm, yshift=3.4mm]current page.south east) {Աղբյուր՝ #1};
  \end{tikzpicture}}

% Speech bubbles (handwritten) and the overlay helper for TikZ:
% \node[visible on=<2->] ... keeps the layout but hides the node until slide 2.
\tikzset{
  say/.style={rectangle callout, draw=plum, line width=0.8pt, fill=cardwhite, rounded corners=3pt,
    inner sep=4pt, align=center, font=\handfont\normalsize, text=plum},
  think/.style={cloud callout, cloud puffs=12, draw=plum, line width=0.8pt, fill=cardwhite,
    inner sep=2pt, align=center, font=\handfont\normalsize, text=plum},
  % boxes for diagrams drawn directly in a frame: a style with #1 cannot be defined inside a frame body
  flow/.style={draw=plum, line width=0.8pt, fill=cardwhite, rounded corners=2mm, align=center,
    font=\small, text=plum, inner sep=4pt, text width=#1},
  flow/.default=3cm,
  chatmsg/.style={draw=plum, line width=0.8pt, rounded corners=3mm, inner sep=6pt, text width=#1,
    align=left, font=\small, text=plum},
  chatmsg/.default=7cm,
  invisible/.style={opacity=0, text opacity=0},
  visible on/.style={alt={#1{}{invisible}}},
  alt/.code args={<#1>#2#3}{\alt<#1>{\pgfkeysalso{#2}}{\pgfkeysalso{#3}}},
}

%  Build from the deck folder:  latexmk -lualatex -outdir=build webinarNN.tex
%  Art (emoji, mascot moods, title banner) comes from common/build_art.py.
%
%  All styling lives here, so decks only hold content. Palette and fonts
%  are the banner's (banner/build_banner.py): plum text on cream, pastel
%  accents, Mardoto for text, Glazz for big titles, Twemoji icons and
%  the pink brain mascot (DECISIONS.md #4-#6).
% =====================================================================

% ---------------------------------------------------------------------
% Fonts (LuaLaTeX + fontspec), shared with the banner
% ---------------------------------------------------------------------
\usepackage{fontspec}
\tracinglostchars=3 % a glyph missing from a font stops the build instead of printing nothing
\defaultfontfeatures{Ligatures=TeX}
\def\idiotfontdir{../../banner/assets/fonts/}
\setsansfont{Mardoto}[Path=\idiotfontdir, Extension=.ttf, UprightFont=*-Regular, BoldFont=*-Bold]
\setmainfont{Mardoto}[Path=\idiotfontdir, Extension=.ttf, UprightFont=*-Regular, BoldFont=*-Bold]
\newfontfamily\mediumfont{Mardoto-Medium}[Path=\idiotfontdir, Extension=.ttf]
\newfontfamily\blackfont{Mardoto-Black}[Path=\idiotfontdir, Extension=.ttf]
\newfontfamily\titlefont{GlazzArm-Regular}[Path=\idiotfontdir, Extension=.ttf] % Armenian letters only
\newfontfamily\handfont{Ink Free} % Windows font with Armenian; speech bubbles only
% Mardoto has no italic: emphasis is medium weight in the accent colour instead.
\renewcommand{\emph}[1]{{\mediumfont\color{plum}#1}}
% No hyphenation: split words are hard to read on a shared screen.
\hyphenpenalty=10000
\exhyphenpenalty=10000

% ---------------------------------------------------------------------
% Packages
% ---------------------------------------------------------------------
\usepackage{graphicx}
\usepackage{booktabs}
\usepackage[most]{tcolorbox}
\usepackage{tikz}
\usetikzlibrary{arrows.meta, positioning, calc, shapes.geometric, shapes.callouts, shapes.misc,
  decorations.pathmorphing, decorations.pathreplacing, backgrounds, fit, patterns.meta}
\usepackage{pgfplots}
\pgfplotsset{compat=1.18}
\usepgfplotslibrary{fillbetween}
\graphicspath{{../common/art/}{images/}}

% ---------------------------------------------------------------------
% Colours: the banner palette. Charts with 3+ series use the Armenian
% flag colours (amRed / amBlue / amOrange).
% ---------------------------------------------------------------------
\definecolor{plum}{HTML}{4A2340}
\definecolor{brainpink}{HTML}{F7A8C4}
\definecolor{fold}{HTML}{E57FA6}
\definecolor{cream}{HTML}{FFF6EC}
\definecolor{cardwhite}{HTML}{FFFDF9}
\definecolor{peach}{HTML}{FFD9C2}
\definecolor{lavender}{HTML}{D9C8F5}
\definecolor{mint}{HTML}{B5E8D5}
\definecolor{butter}{HTML}{FFE8A3}
\definecolor{blush}{HTML}{FF8FA8}
\definecolor{amRed}{HTML}{D90012}
\definecolor{amBlue}{HTML}{0033A0}
\definecolor{amOrange}{HTML}{F2A800}
% darker partners for text on cream
\colorlet{mintdark}{mint!45!black!80!green}
\definecolor{mintdark}{HTML}{2E7D62}
\definecolor{lavenderdark}{HTML}{5B3FA0}
\definecolor{peachdark}{HTML}{B5532A}
\definecolor{pinkdark}{HTML}{C2185B}
% names used by the figure files
\colorlet{ink}{plum}
\colorlet{paper}{cream}
\colorlet{mist}{lavender!40!cream}
\definecolor{night}{HTML}{2B2150}
\colorlet{softgray}{plum!60!cream}
% accent of the current part; \partaccent{<colour>} switches it
\colorlet{accent}{butter}
\newcommand{\partaccent}[1]{\colorlet{accent}{#1}}

% ---------------------------------------------------------------------
% Art: \emoji, \mascot (pages of common/art/*.pdf, see build_art.py)
% ---------------------------------------------------------------------
\makeatletter
\newcommand{\artpage}[3]{\expandafter\def\csname idiot@art@#1@#2\endcsname{#3}}
% generated by build_art.py - do not edit
\artpage{emoji}{brain}{1}
\artpage{emoji}{warning}{2}
\artpage{emoji}{rollercoaster}{3}
\artpage{emoji}{doughnut}{4}
\artpage{emoji}{sleeping}{5}
\artpage{emoji}{scream}{6}
\artpage{emoji}{thought}{7}
\artpage{emoji}{question}{8}
\artpage{emoji}{cyclone}{9}
\artpage{emoji}{calendar}{10}
\artpage{emoji}{clock630}{11}
\artpage{emoji}{computer}{12}
\artpage{emoji}{receipt}{13}
\artpage{emoji}{gift}{14}
\artpage{emoji}{lizard}{15}
\artpage{emoji}{boss}{16}
\artpage{emoji}{keyboard}{17}
\artpage{emoji}{wastebasket}{18}
\artpage{emoji}{batterylow}{19}
\artpage{emoji}{skull}{20}
\artpage{emoji}{car}{21}
\artpage{emoji}{ship}{22}
\artpage{emoji}{airplane}{23}
\artpage{emoji}{nauseated}{24}
\artpage{emoji}{vomit}{25}
\artpage{emoji}{eye}{26}
\artpage{emoji}{ear}{27}
\artpage{emoji}{tree}{28}
\artpage{emoji}{wave}{29}
\artpage{emoji}{book}{30}
\artpage{emoji}{walking}{31}
\artpage{emoji}{dizzy}{32}
\artpage{emoji}{cake}{33}
\artpage{emoji}{icecream}{34}
\artpage{emoji}{potato}{35}
\artpage{emoji}{apple}{36}
\artpage{emoji}{burger}{37}
\artpage{emoji}{salad}{38}
\artpage{emoji}{clock12}{39}
\artpage{emoji}{forkknife}{40}
\artpage{emoji}{pasta}{41}
\artpage{emoji}{fire}{42}
\artpage{emoji}{drool}{43}
\artpage{emoji}{pizza}{44}
\artpage{emoji}{bed}{45}
\artpage{emoji}{personbed}{46}
\artpage{emoji}{moon}{47}
\artpage{emoji}{sun}{48}
\artpage{emoji}{zzz}{49}
\artpage{emoji}{dolphin}{50}
\artpage{emoji}{jellyfish}{51}
\artpage{emoji}{rat}{52}
\artpage{emoji}{chipmunk}{53}
\artpage{emoji}{alarm}{54}
\artpage{emoji}{phone}{55}
\artpage{emoji}{ghost}{56}
\artpage{emoji}{alien}{57}
\artpage{emoji}{yawn}{58}
\artpage{emoji}{baby}{59}
\artpage{emoji}{construction}{60}
\artpage{emoji}{broom}{61}
\artpage{emoji}{thermometer}{62}
\artpage{emoji}{cold}{63}
\artpage{emoji}{monkey}{64}
\artpage{emoji}{tiger}{65}
\artpage{emoji}{axe}{66}
\artpage{emoji}{balloon}{67}
\artpage{emoji}{heart}{68}
\artpage{emoji}{anatomicalheart}{69}
\artpage{emoji}{lungs}{70}
\artpage{emoji}{sweat}{71}
\artpage{emoji}{eyes}{72}
\artpage{emoji}{butterfly}{73}
\artpage{emoji}{runner}{74}
\artpage{emoji}{boxing}{75}
\artpage{emoji}{station}{76}
\artpage{emoji}{taxi}{77}
\artpage{emoji}{bell}{78}
\artpage{emoji}{couch}{79}
\artpage{emoji}{muscle}{80}
\artpage{emoji}{toilet}{81}
\artpage{emoji}{floppy}{82}
\artpage{emoji}{door}{83}
\artpage{emoji}{coffee}{84}
\artpage{emoji}{diving}{85}
\artpage{emoji}{wine}{86}
\artpage{emoji}{beer}{87}
\artpage{emoji}{dove}{88}
\artpage{emoji}{puzzle}{89}
\artpage{emoji}{fish}{90}
\artpage{emoji}{trophy}{91}
\artpage{emoji}{detective}{92}
\artpage{emoji}{judge}{93}
\artpage{emoji}{cheese}{94}
\artpage{emoji}{film}{95}
\artpage{emoji}{magnifier}{96}
\artpage{emoji}{key}{97}
\artpage{emoji}{firework}{98}
\artpage{emoji}{house}{99}
\artpage{emoji}{gloves}{100}
\artpage{emoji}{trafficlight}{101}
\artpage{emoji}{stop}{102}
\artpage{emoji}{mute}{103}
\artpage{emoji}{videogame}{104}
\artpage{emoji}{mirror}{105}
\artpage{emoji}{star}{106}
\artpage{emoji}{hospital}{107}
\artpage{emoji}{pill}{108}
\artpage{emoji}{camera}{109}
\artpage{emoji}{tv}{110}
\artpage{emoji}{crown}{111}
\artpage{emoji}{nerd}{112}
\artpage{emoji}{lock}{113}
\artpage{emoji}{musical}{114}
\artpage{emoji}{chips}{115}
\artpage{emoji}{check}{116}
\artpage{emoji}{cross}{117}
\artpage{emoji}{hand}{118}
\artpage{emoji}{wavehand}{119}
\artpage{emoji}{speech}{120}
\artpage{emoji}{pencil}{121}
\artpage{emoji}{memo}{122}
\artpage{emoji}{monocle}{123}
\artpage{emoji}{thinking}{124}
\artpage{emoji}{grin}{125}
\artpage{emoji}{sweatsmile}{126}
\artpage{emoji}{facepalm}{127}
\artpage{emoji}{popcorn}{128}
\artpage{emoji}{bulb}{129}
\artpage{emoji}{sparkles}{130}
\artpage{emoji}{party}{131}
\artpage{emoji}{chart}{132}
\artpage{emoji}{testtube}{133}
\artpage{emoji}{people}{134}
\artpage{emoji}{pin}{135}
\artpage{emoji}{hourglass}{136}
\artpage{emoji}{microscope}{137}
\artpage{emoji}{books}{138}
\artpage{emoji}{megaphone}{139}
\artpage{emoji}{student}{140}
\artpage{emoji}{zany}{141}
\artpage{emoji}{hundred}{142}
\artpage{emoji}{link}{143}
\artpage{emoji}{target}{144}
\artpage{emoji}{clover}{145}
\artpage{emoji}{ufo}{146}
\artpage{emoji}{spider}{147}
\artpage{emoji}{spiderweb}{148}
\artpage{emoji}{snake}{149}
\artpage{emoji}{microphone}{150}
\artpage{emoji}{clown}{151}
\artpage{emoji}{pumpkin}{152}
\artpage{emoji}{knife}{153}
\artpage{emoji}{bat}{154}
\artpage{emoji}{blackcat}{155}
\artpage{emoji}{crossedfingers}{156}
\artpage{emoji}{ladder}{157}
\artpage{emoji}{thumbsup}{158}
\artpage{emoji}{thumbsdown}{159}
\artpage{emoji}{rocket}{160}
\artpage{emoji}{satellite}{161}
\artpage{emoji}{newspaper}{162}
\artpage{emoji}{chocolate}{163}
\artpage{emoji}{syringe}{164}
\artpage{emoji}{dice}{165}
\artpage{emoji}{hauntedhouse}{166}
\artpage{emoji}{flushed}{167}
\artpage{emoji}{fearful}{168}
\artpage{emoji}{graduation}{169}
\artpage{emoji}{teacher}{170}
\artpage{emoji}{scientist}{171}
\artpage{emoji}{chess}{172}
\artpage{emoji}{abacus}{173}
\artpage{emoji}{ruler}{174}
\artpage{emoji}{jeans}{175}
\artpage{emoji}{glasses}{176}
\artpage{emoji}{dna}{177}
\artpage{emoji}{loudspeaker}{178}
\artpage{emoji}{speakinghead}{179}
\artpage{emoji}{seedling}{180}
\artpage{emoji}{weightlifter}{181}
\artpage{emoji}{giraffe}{182}
\artpage{emoji}{mouse}{183}
\artpage{emoji}{exploding}{184}
\artpage{emoji}{owl}{185}
\artpage{emoji}{clipboard}{186}
\artpage{emoji}{numbers}{187}
\artpage{emoji}{medal}{188}
\artpage{emoji}{rose}{189}
\artpage{emoji}{nose}{190}
\artpage{emoji}{tongue}{191}
\artpage{emoji}{yum}{192}
\artpage{emoji}{bread}{193}
\artpage{emoji}{lemon}{194}
\artpage{emoji}{chili}{195}
\artpage{emoji}{salt}{196}
\artpage{emoji}{candy}{197}
\artpage{emoji}{speaker}{198}
\artpage{emoji}{headphones}{199}
\artpage{emoji}{feather}{200}
\artpage{emoji}{gorilla}{201}
\artpage{emoji}{basketball}{202}
\artpage{emoji}{magicwand}{203}
\artpage{emoji}{tophat}{204}
\artpage{emoji}{cloud}{205}
\artpage{emoji}{plug}{206}
\artpage{emoji}{clinking}{207}
\artpage{emoji}{dog}{208}
\artpage{emoji}{rainbow}{209}
\artpage{emoji}{slightsmile}{210}
\artpage{emoji}{palette}{211}
\artpage{emoji}{masks}{212}
\artpage{emoji}{angry}{213}
\artpage{emoji}{rage}{214}
\artpage{emoji}{steam}{215}
\artpage{emoji}{bomb}{216}
\artpage{emoji}{volcano}{217}
\artpage{emoji}{angersymbol}{218}
\artpage{emoji}{crystalball}{219}
\artpage{emoji}{railway}{220}
\artpage{emoji}{worker}{221}
\artpage{emoji}{laughing}{222}
\artpage{emoji}{rofl}{223}
\artpage{emoji}{banana}{224}
\artpage{emoji}{mountain}{225}
\artpage{emoji}{moneybag}{226}
\artpage{emoji}{lightning}{227}
\artpage{emoji}{joker}{228}
\artpage{emoji}{finishflag}{229}
\artpage{emoji}{hug}{230}
\artpage{emoji}{handshake}{231}
\artpage{emoji}{lying}{232}
\artpage{emoji}{neutral}{233}
\artpage{emoji}{carrot}{234}
\artpage{emoji}{wood}{235}
\artpage{emoji}{brokenheart}{236}
\artpage{emoji}{loveletter}{237}
\artpage{emoji}{couple}{238}
\artpage{emoji}{kiss}{239}
\artpage{emoji}{cowboy}{240}
\artpage{emoji}{guitar}{241}
\artpage{emoji}{family}{242}
\artpage{emoji}{police}{243}
\artpage{emoji}{sheep}{244}
\artpage{emoji}{clapping}{245}
\artpage{emoji}{imp}{246}
\artpage{emoji}{ballot}{247}
\artpage{emoji}{smileheart}{248}
\artpage{emoji}{dollar}{249}
\artpage{emoji}{earth}{250}
\artpage{emoji}{raincloud}{251}
\artpage{emoji}{sunbehindcloud}{252}
\artpage{emoji}{pensive}{253}
\artpage{emoji}{cry}{254}
\artpage{emoji}{bandage}{255}
\artpage{emoji}{stethoscope}{256}
\artpage{emoji}{cigarette}{257}
\artpage{emoji}{slotmachine}{258}
\artpage{emoji}{siren}{259}
\artpage{emoji}{sunrise}{260}
\artpage{emoji}{healthworker}{261}
\artpage{emoji}{hearnoevil}{262}
\artpage{emoji}{anchor}{263}
\artpage{emoji}{candle}{264}
\artpage{emoji}{telephone}{265}
\artpage{emoji}{battery}{266}
\artpage{emoji}{weary}{267}
\artpage{emoji}{lotus}{268}
\artpage{mascot}{panic}{1}
\artpage{mascot}{happy}{2}
\artpage{mascot}{scared}{3}
\artpage{mascot}{sleepy}{4}
\artpage{mascot}{sick}{5}
\artpage{mascot}{confused}{6}
\artpage{mascot}{smug}{7}
\artpage{mascot}{hungry}{8}
\artpage{mascot}{thinking}{9}
\artpage{mascot}{sad}{10}
\artpage{mascot}{angry}{11}

\newcommand{\idiot@art}[3]{% kind, name, size options
  \ifcsname idiot@art@#1@#2\endcsname
    \includegraphics[#3, page=\csname idiot@art@#1@#2\endcsname]{#1.pdf}%
  \else
    \PackageError{idiot}{Unknown #1 '#2'; add it to build_art.py and rebuild the art}{}%
  \fi}
% \emoji[<height>]{name}: inline icon, sits on the text baseline like a letter
\newcommand{\emoji}[2][1.15em]{\raisebox{-0.2em}{\idiot@art{emoji}{#2}{height=#1}}}
% \bigemoji[<height>]{name}: no baseline shift, for TikZ nodes and big pictures
\newcommand{\bigemoji}[2][1.5cm]{\idiot@art{emoji}{#2}{height=#1}}
% \mascot[<width>]{mood}: panic, happy, scared, sleepy, sick, confused, smug, hungry, thinking
\newcommand{\mascot}[2][2.6cm]{\idiot@art{mascot}{#2}{width=#1}}
\makeatother
% the older figure files place the mascot with \pic{noggin=<mood>}
\tikzset{pics/noggin/.style={code={\node[inner sep=0, transform shape] {\mascot[2.7cm]{#1}};}}}

% ---------------------------------------------------------------------
% Beamer look
% ---------------------------------------------------------------------
\setbeamertemplate{navigation symbols}{}
\setbeamercolor{normal text}{fg=plum, bg=cream}
\setbeamercolor{structure}{fg=plum}
\setbeamercolor{alerted text}{fg=pinkdark}
\setbeamerfont{alerted text}{series=\bfseries}
\setbeamercolor{frametitle}{fg=plum}
\setbeamersize{text margin left=9mm, text margin right=9mm}
\setbeamertemplate{frame begin}{\raggedright} % justified Armenian leaves big gaps

% cream paper with the banner's faint dot grid
\setbeamertemplate{background}{%
  \begin{tikzpicture}
    \fill[cream] (0,0) rectangle (\paperwidth,\paperheight);
    \fill[pattern={Dots[distance=4.4mm, radius=0.28mm]}, pattern color=plum!9!cream]
      (0,0) rectangle (\paperwidth,\paperheight);
  \end{tikzpicture}}

% topic pill, top right of every frame: set by \topic (emoji + short name)
\def\idiottopicemoji{}
\def\idiottopicname{}
\newcommand{\idiotpill}[3][cardwhite]{% fill, emoji, text
  \tikz[baseline=(p.base)]{\node[fill=#1, draw=plum, line width=0.7pt, rounded corners=2.6mm,
      inner xsep=2.2mm, inner ysep=1.1mm, font=\mediumfont\footnotesize, text=plum] (p)
      {\ifx\relax#2\relax\else\emoji[1.05em]{#2}\hspace{1.2mm}\fi#3};
    \begin{scope}[on background layer]
      \fill[plum, rounded corners=2.6mm] ([xshift=0.7mm, yshift=-0.7mm]p.south west) rectangle ([xshift=0.7mm, yshift=-0.7mm]p.north east);
    \end{scope}}}
% \chip{emoji}{text}: inline sticker label, same as the banner's topic chips
\newcommand{\chip}[3][cardwhite]{\idiotpill[#1]{#2}{#3}}

\setbeamertemplate{frametitle}{%
  \vskip 3.5mm
  \begin{minipage}[t]{0.66\paperwidth}
    \raggedright{\blackfont\fontsize{17.5}{20.5}\selectfont\insertframetitle}\par
    \vskip 1.3mm
    \tikz\fill[accent, rounded corners=0.6mm] (0,0) rectangle (1.5cm,1.1mm);
  \end{minipage}%
  \hfill
  \ifx\idiottopicname\empty\else
    \raisebox{-1mm}{\idiotpill{\idiottopicemoji}{\idiottopicname}}%
  \fi
  \vskip 1mm}

\setbeamertemplate{footline}{%
  \hbox to \paperwidth{\hskip 9mm
    {\mediumfont\fontsize{7}{8}\selectfont\color{plum!55!cream}Ապուշ ուղեղը · \insertshorttitle}\hfill
    {\blackfont\fontsize{8}{9}\selectfont\color{plum!70!cream}\insertframenumber}\hskip 9mm}
  \vskip 3mm}

\setbeamertemplate{itemize item}{\tikz[baseline=-0.65ex]\draw[plum, fill=accent, line width=0.6pt] (0,0) circle (0.95mm);}
\setbeamertemplate{itemize subitem}{\tikz[baseline=-0.65ex]\fill[plum!70] (0,0) circle (0.6mm);}
\setbeamertemplate{enumerate item}{\tikz[baseline=(n.base)]\node[fill=accent, draw=plum, line width=0.6pt,
  circle, inner sep=0.5mm, minimum size=4.2mm, font=\blackfont\scriptsize] (n) {\insertenumlabel};}
\setlength{\leftmargini}{5mm}
\setbeamertemplate{itemize/enumerate body begin}{\setlength{\itemsep}{1.6mm}}
\setbeamertemplate{itemize/enumerate subbody begin}{\setlength{\itemsep}{0.6mm}\small}

% ---------------------------------------------------------------------
% Structure: parts (full-colour openers) and topics (book sections)
% ---------------------------------------------------------------------
% \topic{<emoji>}{<short name>}: starts a book section; the pill shows on every following frame
\newcommand{\topic}[2]{\subsection{#2}\gdef\idiottopicemoji{#1}\gdef\idiottopicname{#2}}
\newcommand{\notopic}{\gdef\idiottopicemoji{}\gdef\idiottopicname{}}

% \partopener{<label>}{<title>}{<subtitle>}{<mood>}{<colour>}
\newcommand{\partopener}[5]{%
  \section{#2}\notopic\partaccent{#5}%
  \begin{frame}[plain]
    \begin{tikzpicture}[remember picture, overlay]
      \fill[#5] (current page.south west) rectangle (current page.north east);
      \fill[pattern={Dots[distance=4.4mm, radius=0.28mm]}, pattern color=plum!12!#5]
        (current page.south west) rectangle (current page.north east);
      \fill[cardwhite, opacity=0.45] ([xshift=-3.3cm, yshift=-0.2cm]current page.east) circle (3.1cm);
      \node[anchor=north west, inner sep=0] at ([xshift=10mm, yshift=-12mm]current page.north west)
        {\idiotpill{}{#1}};
      \node[anchor=west, text=plum, align=left, text width=9.2cm, inner sep=0]
        (t) at ([xshift=10mm, yshift=5mm]current page.west)
        {\titlefont\fontsize{38}{40}\selectfont #2};
      \node[anchor=north west, text=plum, text width=9.2cm, inner sep=0, align=left]
        at ([yshift=-5mm]t.south west) {\mediumfont\large #3};
      \node at ([xshift=-3.3cm, yshift=-0.2cm]current page.east) {\mascot[4.6cm]{#4}};
    \end{tikzpicture}
  \end{frame}}

% \topicopener{<emoji>}{<short name>}{<book's section title>}{<what it is about>}{<mood>}
\newcommand{\topicopener}[5]{%
  \topic{#1}{#2}%
  \begin{frame}[plain]
    \begin{tikzpicture}[remember picture, overlay]
      \fill[accent!55!cream] (current page.south west) rectangle (current page.north east);
      \fill[pattern={Dots[distance=4.4mm, radius=0.28mm]}, pattern color=plum!10!accent!55!cream]
        (current page.south west) rectangle (current page.north east);
      \node[anchor=north west, inner sep=0] at ([xshift=10mm, yshift=-12mm]current page.north west)
        {\idiotpill{#1}{#2}};
      \node[anchor=west, text=plum, text width=9.6cm, align=left, inner sep=0]
        (t) at ([xshift=10mm, yshift=6mm]current page.west)
        {\blackfont\fontsize{24}{28}\selectfont «#3»};
      \node[anchor=north west, text=plum, text width=9.6cm, align=left, inner sep=0]
        at ([yshift=-4mm]t.south west) {\mediumfont\large #4};
      \node at ([xshift=-3.0cm, yshift=-0.4cm]current page.east) {\mascot[4.2cm]{#5}};
      \node[rotate=-12] at ([xshift=-1.4cm, yshift=2.3cm]current page.east) {\bigemoji[1.5cm]{#1}};
    \end{tikzpicture}
  \end{frame}}

% ---------------------------------------------------------------------
% Boxes (sticker style: plum outline + hard plum shadow, like the banner)
% ---------------------------------------------------------------------
\tcbset{
  sticker/.style={enhanced, arc=3mm, boxrule=0.9pt, colframe=plum, colback=cardwhite, halign=left,
    left=3mm, right=3mm, top=2mm, bottom=2mm, fontupper=\normalsize, coltext=plum,
    shadow={1.1mm}{-1.1mm}{0mm}{fill=plum}},
  tagged/.style={sticker, colback=#1!45!cream, colbacktitle=#1, coltitle=plum,
    fonttitle=\blackfont\footnotesize, top=3mm,
    attach boxed title to top left={xshift=4mm, yshift*=-2.4mm},
    boxed title style={arc=2.2mm, boxrule=0.9pt, colframe=plum, left=1.5mm, right=1.5mm, top=0.6mm, bottom=0.6mm}},
}
\newtcolorbox{card}[1][]{sticker, #1}
\newtcolorbox{tryit}[1][]{tagged=mint, title={\emoji[1em]{testtube}\ Փորձիր հիմա}, #1}
\newtcolorbox{poll}[1][]{tagged=butter, title={\emoji[1em]{hand}\ Հարցում}, #1}
\newtcolorbox{chatbox}[1][]{tagged=lavender, title={\emoji[1em]{speech}\ Գրիր չաթում}, #1}
\newtcolorbox{scienceupdate}[1][]{tagged=peach, fontupper=\small, title={\emoji[1em]{microscope}\ Գիտությունն այսօր}, #1}
\newtcolorbox{bookquote}[1][]{sticker, colback=cardwhite, fontupper=\mediumfont, left=4mm,
  borderline west={1.4mm}{0pt}{accent},
  before upper={{\blackfont\footnotesize\emoji[1em]{book}\ Բըրնեթը գրում է}\par\smallskip}, #1}

% Brain.exe-style popup window: \begin{popup}{<window title>} ... \end{popup}
\newcommand{\winbuttons}{%
  \tikz[baseline=-0.6ex, line width=0.6pt]{
    \foreach \i/\f in {0/cardwhite, 1/cardwhite, 2/blush}{
      \filldraw[fill=\f, draw=plum, rounded corners=0.6mm] (\i*4.3mm,-1.6mm) rectangle ++(3.3mm,3.3mm);}
    \draw[plum] (0.8mm,-0.8mm) -- (2.5mm,-0.8mm);
    \draw[plum] (5.1mm,-0.8mm) rectangle (6.8mm,0.9mm);
    \draw[plum] (9.4mm,-0.9mm) -- (11.1mm,0.8mm) (9.4mm,0.8mm) -- (11.1mm,-0.9mm);}}
\newtcolorbox{popup}[2][]{sticker, colback=cardwhite, colbacktitle=brainpink, coltitle=plum,
  fonttitle=\blackfont\small, title={\emoji[1.1em]{warning}\hspace{1.5mm}#2\hfill\winbuttons},
  titlerule=0.9pt, toptitle=0.8mm, bottomtitle=0.8mm, #1}

% Myth or fact: \mythorfact{<statement>}{<verdict>}{<colour>}{<explanation>}; verdict shows on slide 2
\newcommand{\mythorfact}[4]{%
  \begin{center}
  \begin{tikzpicture}
    \node[draw=plum, line width=0.9pt, fill=cardwhite, rounded corners=4mm, inner xsep=8mm, inner ysep=6mm,
          text width=10cm, align=center, font=\blackfont\Large, text=plum] (mf) {«#1»};
    \begin{scope}[on background layer]
      \fill[plum, rounded corners=4mm] ([xshift=1.2mm, yshift=-1.2mm]mf.south west) rectangle ([xshift=1.2mm, yshift=-1.2mm]mf.north east);
    \end{scope}
    \node[anchor=west] at ([xshift=6mm]mf.north west) {\idiotpill[butter]{question}{Առասպել, թե՞ փաստ}};
    \onslide<2->{\node[draw=#3, text=#3, line width=2pt, rounded corners=2mm, rotate=-10, font=\blackfont\LARGE,
      inner sep=4pt, fill=cardwhite] at ([xshift=-8mm, yshift=1mm]mf.south east) {#2};}
  \end{tikzpicture}
  \end{center}
  \onslide<2->{#4}}

% Big question with the confused mascot: \bigquestion{<question>}{<small line>}
\newcommand{\bigquestion}[2]{%
  \vfill
  \begin{center}
  \begin{tikzpicture}
    \node[text width=9.6cm, align=center, font=\blackfont\fontsize{24}{29}\selectfont, text=plum] (q) {#1};
    \node[below=4mm of q, text width=9.6cm, align=center, font=\mediumfont\large, text=softgray] {#2};
    \node[anchor=east] at ([xshift=-2mm]q.west) {\mascot[2.6cm]{confused}};
  \end{tikzpicture}
  \end{center}
  \vfill}

% Small source line at the bottom of a frame: \source{Loftus & Palmer, 1974}
\newcommand{\source}[1]{%
  \begin{tikzpicture}[remember picture, overlay]
    \node[anchor=south east, font=\fontsize{6.5}{7.5}\selectfont, text=plum!60!cream, inner sep=0]
      at ([xshift=-17mm, yshift=3.4mm]current page.south east) {Աղբյուր՝ #1};
  \end{tikzpicture}}

% Speech bubbles (handwritten) and the overlay helper for TikZ:
% \node[visible on=<2->] ... keeps the layout but hides the node until slide 2.
\tikzset{
  say/.style={rectangle callout, draw=plum, line width=0.8pt, fill=cardwhite, rounded corners=3pt,
    inner sep=4pt, align=center, font=\handfont\normalsize, text=plum},
  think/.style={cloud callout, cloud puffs=12, draw=plum, line width=0.8pt, fill=cardwhite,
    inner sep=2pt, align=center, font=\handfont\normalsize, text=plum},
  % boxes for diagrams drawn directly in a frame: a style with #1 cannot be defined inside a frame body
  flow/.style={draw=plum, line width=0.8pt, fill=cardwhite, rounded corners=2mm, align=center,
    font=\small, text=plum, inner sep=4pt, text width=#1},
  flow/.default=3cm,
  chatmsg/.style={draw=plum, line width=0.8pt, rounded corners=3mm, inner sep=6pt, text width=#1,
    align=left, font=\small, text=plum},
  chatmsg/.default=7cm,
  invisible/.style={opacity=0, text opacity=0},
  visible on/.style={alt={#1{}{invisible}}},
  alt/.code args={<#1>#2#3}{\alt<#1>{\pgfkeysalso{#2}}{\pgfkeysalso{#3}}},
}

%  Build from the deck folder:  latexmk -lualatex -outdir=build webinarNN.tex
%  Art (emoji, mascot moods, title banner) comes from common/build_art.py.
%
%  All styling lives here, so decks only hold content. Palette and fonts
%  are the banner's (banner/build_banner.py): plum text on cream, pastel
%  accents, Mardoto for text, Glazz for big titles, Twemoji icons and
%  the pink brain mascot (DECISIONS.md #4-#6).
% =====================================================================

% ---------------------------------------------------------------------
% Fonts (LuaLaTeX + fontspec), shared with the banner
% ---------------------------------------------------------------------
\usepackage{fontspec}
\tracinglostchars=3 % a glyph missing from a font stops the build instead of printing nothing
\defaultfontfeatures{Ligatures=TeX}
\def\idiotfontdir{../../banner/assets/fonts/}
\setsansfont{Mardoto}[Path=\idiotfontdir, Extension=.ttf, UprightFont=*-Regular, BoldFont=*-Bold]
\setmainfont{Mardoto}[Path=\idiotfontdir, Extension=.ttf, UprightFont=*-Regular, BoldFont=*-Bold]
\newfontfamily\mediumfont{Mardoto-Medium}[Path=\idiotfontdir, Extension=.ttf]
\newfontfamily\blackfont{Mardoto-Black}[Path=\idiotfontdir, Extension=.ttf]
\newfontfamily\titlefont{GlazzArm-Regular}[Path=\idiotfontdir, Extension=.ttf] % Armenian letters only
\newfontfamily\handfont{Ink Free} % Windows font with Armenian; speech bubbles only
% Mardoto has no italic: emphasis is medium weight in the accent colour instead.
\renewcommand{\emph}[1]{{\mediumfont\color{plum}#1}}
% No hyphenation: split words are hard to read on a shared screen.
\hyphenpenalty=10000
\exhyphenpenalty=10000

% ---------------------------------------------------------------------
% Packages
% ---------------------------------------------------------------------
\usepackage{graphicx}
\usepackage{booktabs}
\usepackage[most]{tcolorbox}
\usepackage{tikz}
\usetikzlibrary{arrows.meta, positioning, calc, shapes.geometric, shapes.callouts, shapes.misc,
  decorations.pathmorphing, decorations.pathreplacing, backgrounds, fit, patterns.meta}
\usepackage{pgfplots}
\pgfplotsset{compat=1.18}
\usepgfplotslibrary{fillbetween}
\graphicspath{{../common/art/}{images/}}

% ---------------------------------------------------------------------
% Colours: the banner palette. Charts with 3+ series use the Armenian
% flag colours (amRed / amBlue / amOrange).
% ---------------------------------------------------------------------
\definecolor{plum}{HTML}{4A2340}
\definecolor{brainpink}{HTML}{F7A8C4}
\definecolor{fold}{HTML}{E57FA6}
\definecolor{cream}{HTML}{FFF6EC}
\definecolor{cardwhite}{HTML}{FFFDF9}
\definecolor{peach}{HTML}{FFD9C2}
\definecolor{lavender}{HTML}{D9C8F5}
\definecolor{mint}{HTML}{B5E8D5}
\definecolor{butter}{HTML}{FFE8A3}
\definecolor{blush}{HTML}{FF8FA8}
\definecolor{amRed}{HTML}{D90012}
\definecolor{amBlue}{HTML}{0033A0}
\definecolor{amOrange}{HTML}{F2A800}
% darker partners for text on cream
\colorlet{mintdark}{mint!45!black!80!green}
\definecolor{mintdark}{HTML}{2E7D62}
\definecolor{lavenderdark}{HTML}{5B3FA0}
\definecolor{peachdark}{HTML}{B5532A}
\definecolor{pinkdark}{HTML}{C2185B}
% names used by the figure files
\colorlet{ink}{plum}
\colorlet{paper}{cream}
\colorlet{mist}{lavender!40!cream}
\definecolor{night}{HTML}{2B2150}
\colorlet{softgray}{plum!60!cream}
% accent of the current part; \partaccent{<colour>} switches it
\colorlet{accent}{butter}
\newcommand{\partaccent}[1]{\colorlet{accent}{#1}}

% ---------------------------------------------------------------------
% Art: \emoji, \mascot (pages of common/art/*.pdf, see build_art.py)
% ---------------------------------------------------------------------
\makeatletter
\newcommand{\artpage}[3]{\expandafter\def\csname idiot@art@#1@#2\endcsname{#3}}
% generated by build_art.py - do not edit
\artpage{emoji}{brain}{1}
\artpage{emoji}{warning}{2}
\artpage{emoji}{rollercoaster}{3}
\artpage{emoji}{doughnut}{4}
\artpage{emoji}{sleeping}{5}
\artpage{emoji}{scream}{6}
\artpage{emoji}{thought}{7}
\artpage{emoji}{question}{8}
\artpage{emoji}{cyclone}{9}
\artpage{emoji}{calendar}{10}
\artpage{emoji}{clock630}{11}
\artpage{emoji}{computer}{12}
\artpage{emoji}{receipt}{13}
\artpage{emoji}{gift}{14}
\artpage{emoji}{lizard}{15}
\artpage{emoji}{boss}{16}
\artpage{emoji}{keyboard}{17}
\artpage{emoji}{wastebasket}{18}
\artpage{emoji}{batterylow}{19}
\artpage{emoji}{skull}{20}
\artpage{emoji}{car}{21}
\artpage{emoji}{ship}{22}
\artpage{emoji}{airplane}{23}
\artpage{emoji}{nauseated}{24}
\artpage{emoji}{vomit}{25}
\artpage{emoji}{eye}{26}
\artpage{emoji}{ear}{27}
\artpage{emoji}{tree}{28}
\artpage{emoji}{wave}{29}
\artpage{emoji}{book}{30}
\artpage{emoji}{walking}{31}
\artpage{emoji}{dizzy}{32}
\artpage{emoji}{cake}{33}
\artpage{emoji}{icecream}{34}
\artpage{emoji}{potato}{35}
\artpage{emoji}{apple}{36}
\artpage{emoji}{burger}{37}
\artpage{emoji}{salad}{38}
\artpage{emoji}{clock12}{39}
\artpage{emoji}{forkknife}{40}
\artpage{emoji}{pasta}{41}
\artpage{emoji}{fire}{42}
\artpage{emoji}{drool}{43}
\artpage{emoji}{pizza}{44}
\artpage{emoji}{bed}{45}
\artpage{emoji}{personbed}{46}
\artpage{emoji}{moon}{47}
\artpage{emoji}{sun}{48}
\artpage{emoji}{zzz}{49}
\artpage{emoji}{dolphin}{50}
\artpage{emoji}{jellyfish}{51}
\artpage{emoji}{rat}{52}
\artpage{emoji}{chipmunk}{53}
\artpage{emoji}{alarm}{54}
\artpage{emoji}{phone}{55}
\artpage{emoji}{ghost}{56}
\artpage{emoji}{alien}{57}
\artpage{emoji}{yawn}{58}
\artpage{emoji}{baby}{59}
\artpage{emoji}{construction}{60}
\artpage{emoji}{broom}{61}
\artpage{emoji}{thermometer}{62}
\artpage{emoji}{cold}{63}
\artpage{emoji}{monkey}{64}
\artpage{emoji}{tiger}{65}
\artpage{emoji}{axe}{66}
\artpage{emoji}{balloon}{67}
\artpage{emoji}{heart}{68}
\artpage{emoji}{anatomicalheart}{69}
\artpage{emoji}{lungs}{70}
\artpage{emoji}{sweat}{71}
\artpage{emoji}{eyes}{72}
\artpage{emoji}{butterfly}{73}
\artpage{emoji}{runner}{74}
\artpage{emoji}{boxing}{75}
\artpage{emoji}{station}{76}
\artpage{emoji}{taxi}{77}
\artpage{emoji}{bell}{78}
\artpage{emoji}{couch}{79}
\artpage{emoji}{muscle}{80}
\artpage{emoji}{toilet}{81}
\artpage{emoji}{floppy}{82}
\artpage{emoji}{door}{83}
\artpage{emoji}{coffee}{84}
\artpage{emoji}{diving}{85}
\artpage{emoji}{wine}{86}
\artpage{emoji}{beer}{87}
\artpage{emoji}{dove}{88}
\artpage{emoji}{puzzle}{89}
\artpage{emoji}{fish}{90}
\artpage{emoji}{trophy}{91}
\artpage{emoji}{detective}{92}
\artpage{emoji}{judge}{93}
\artpage{emoji}{cheese}{94}
\artpage{emoji}{film}{95}
\artpage{emoji}{magnifier}{96}
\artpage{emoji}{key}{97}
\artpage{emoji}{firework}{98}
\artpage{emoji}{house}{99}
\artpage{emoji}{gloves}{100}
\artpage{emoji}{trafficlight}{101}
\artpage{emoji}{stop}{102}
\artpage{emoji}{mute}{103}
\artpage{emoji}{videogame}{104}
\artpage{emoji}{mirror}{105}
\artpage{emoji}{star}{106}
\artpage{emoji}{hospital}{107}
\artpage{emoji}{pill}{108}
\artpage{emoji}{camera}{109}
\artpage{emoji}{tv}{110}
\artpage{emoji}{crown}{111}
\artpage{emoji}{nerd}{112}
\artpage{emoji}{lock}{113}
\artpage{emoji}{musical}{114}
\artpage{emoji}{chips}{115}
\artpage{emoji}{check}{116}
\artpage{emoji}{cross}{117}
\artpage{emoji}{hand}{118}
\artpage{emoji}{wavehand}{119}
\artpage{emoji}{speech}{120}
\artpage{emoji}{pencil}{121}
\artpage{emoji}{memo}{122}
\artpage{emoji}{monocle}{123}
\artpage{emoji}{thinking}{124}
\artpage{emoji}{grin}{125}
\artpage{emoji}{sweatsmile}{126}
\artpage{emoji}{facepalm}{127}
\artpage{emoji}{popcorn}{128}
\artpage{emoji}{bulb}{129}
\artpage{emoji}{sparkles}{130}
\artpage{emoji}{party}{131}
\artpage{emoji}{chart}{132}
\artpage{emoji}{testtube}{133}
\artpage{emoji}{people}{134}
\artpage{emoji}{pin}{135}
\artpage{emoji}{hourglass}{136}
\artpage{emoji}{microscope}{137}
\artpage{emoji}{books}{138}
\artpage{emoji}{megaphone}{139}
\artpage{emoji}{student}{140}
\artpage{emoji}{zany}{141}
\artpage{emoji}{hundred}{142}
\artpage{emoji}{link}{143}
\artpage{emoji}{target}{144}
\artpage{emoji}{clover}{145}
\artpage{emoji}{ufo}{146}
\artpage{emoji}{spider}{147}
\artpage{emoji}{spiderweb}{148}
\artpage{emoji}{snake}{149}
\artpage{emoji}{microphone}{150}
\artpage{emoji}{clown}{151}
\artpage{emoji}{pumpkin}{152}
\artpage{emoji}{knife}{153}
\artpage{emoji}{bat}{154}
\artpage{emoji}{blackcat}{155}
\artpage{emoji}{crossedfingers}{156}
\artpage{emoji}{ladder}{157}
\artpage{emoji}{thumbsup}{158}
\artpage{emoji}{thumbsdown}{159}
\artpage{emoji}{rocket}{160}
\artpage{emoji}{satellite}{161}
\artpage{emoji}{newspaper}{162}
\artpage{emoji}{chocolate}{163}
\artpage{emoji}{syringe}{164}
\artpage{emoji}{dice}{165}
\artpage{emoji}{hauntedhouse}{166}
\artpage{emoji}{flushed}{167}
\artpage{emoji}{fearful}{168}
\artpage{emoji}{graduation}{169}
\artpage{emoji}{teacher}{170}
\artpage{emoji}{scientist}{171}
\artpage{emoji}{chess}{172}
\artpage{emoji}{abacus}{173}
\artpage{emoji}{ruler}{174}
\artpage{emoji}{jeans}{175}
\artpage{emoji}{glasses}{176}
\artpage{emoji}{dna}{177}
\artpage{emoji}{loudspeaker}{178}
\artpage{emoji}{speakinghead}{179}
\artpage{emoji}{seedling}{180}
\artpage{emoji}{weightlifter}{181}
\artpage{emoji}{giraffe}{182}
\artpage{emoji}{mouse}{183}
\artpage{emoji}{exploding}{184}
\artpage{emoji}{owl}{185}
\artpage{emoji}{clipboard}{186}
\artpage{emoji}{numbers}{187}
\artpage{emoji}{medal}{188}
\artpage{emoji}{rose}{189}
\artpage{emoji}{nose}{190}
\artpage{emoji}{tongue}{191}
\artpage{emoji}{yum}{192}
\artpage{emoji}{bread}{193}
\artpage{emoji}{lemon}{194}
\artpage{emoji}{chili}{195}
\artpage{emoji}{salt}{196}
\artpage{emoji}{candy}{197}
\artpage{emoji}{speaker}{198}
\artpage{emoji}{headphones}{199}
\artpage{emoji}{feather}{200}
\artpage{emoji}{gorilla}{201}
\artpage{emoji}{basketball}{202}
\artpage{emoji}{magicwand}{203}
\artpage{emoji}{tophat}{204}
\artpage{emoji}{cloud}{205}
\artpage{emoji}{plug}{206}
\artpage{emoji}{clinking}{207}
\artpage{emoji}{dog}{208}
\artpage{emoji}{rainbow}{209}
\artpage{emoji}{slightsmile}{210}
\artpage{emoji}{palette}{211}
\artpage{emoji}{masks}{212}
\artpage{emoji}{angry}{213}
\artpage{emoji}{rage}{214}
\artpage{emoji}{steam}{215}
\artpage{emoji}{bomb}{216}
\artpage{emoji}{volcano}{217}
\artpage{emoji}{angersymbol}{218}
\artpage{emoji}{crystalball}{219}
\artpage{emoji}{railway}{220}
\artpage{emoji}{worker}{221}
\artpage{emoji}{laughing}{222}
\artpage{emoji}{rofl}{223}
\artpage{emoji}{banana}{224}
\artpage{emoji}{mountain}{225}
\artpage{emoji}{moneybag}{226}
\artpage{emoji}{lightning}{227}
\artpage{emoji}{joker}{228}
\artpage{emoji}{finishflag}{229}
\artpage{emoji}{hug}{230}
\artpage{emoji}{handshake}{231}
\artpage{emoji}{lying}{232}
\artpage{emoji}{neutral}{233}
\artpage{emoji}{carrot}{234}
\artpage{emoji}{wood}{235}
\artpage{emoji}{brokenheart}{236}
\artpage{emoji}{loveletter}{237}
\artpage{emoji}{couple}{238}
\artpage{emoji}{kiss}{239}
\artpage{emoji}{cowboy}{240}
\artpage{emoji}{guitar}{241}
\artpage{emoji}{family}{242}
\artpage{emoji}{police}{243}
\artpage{emoji}{sheep}{244}
\artpage{emoji}{clapping}{245}
\artpage{emoji}{imp}{246}
\artpage{emoji}{ballot}{247}
\artpage{emoji}{smileheart}{248}
\artpage{emoji}{dollar}{249}
\artpage{emoji}{earth}{250}
\artpage{emoji}{raincloud}{251}
\artpage{emoji}{sunbehindcloud}{252}
\artpage{emoji}{pensive}{253}
\artpage{emoji}{cry}{254}
\artpage{emoji}{bandage}{255}
\artpage{emoji}{stethoscope}{256}
\artpage{emoji}{cigarette}{257}
\artpage{emoji}{slotmachine}{258}
\artpage{emoji}{siren}{259}
\artpage{emoji}{sunrise}{260}
\artpage{emoji}{healthworker}{261}
\artpage{emoji}{hearnoevil}{262}
\artpage{emoji}{anchor}{263}
\artpage{emoji}{candle}{264}
\artpage{emoji}{telephone}{265}
\artpage{emoji}{battery}{266}
\artpage{emoji}{weary}{267}
\artpage{emoji}{lotus}{268}
\artpage{mascot}{panic}{1}
\artpage{mascot}{happy}{2}
\artpage{mascot}{scared}{3}
\artpage{mascot}{sleepy}{4}
\artpage{mascot}{sick}{5}
\artpage{mascot}{confused}{6}
\artpage{mascot}{smug}{7}
\artpage{mascot}{hungry}{8}
\artpage{mascot}{thinking}{9}
\artpage{mascot}{sad}{10}
\artpage{mascot}{angry}{11}

\newcommand{\idiot@art}[3]{% kind, name, size options
  \ifcsname idiot@art@#1@#2\endcsname
    \includegraphics[#3, page=\csname idiot@art@#1@#2\endcsname]{#1.pdf}%
  \else
    \PackageError{idiot}{Unknown #1 '#2'; add it to build_art.py and rebuild the art}{}%
  \fi}
% \emoji[<height>]{name}: inline icon, sits on the text baseline like a letter
\newcommand{\emoji}[2][1.15em]{\raisebox{-0.2em}{\idiot@art{emoji}{#2}{height=#1}}}
% \bigemoji[<height>]{name}: no baseline shift, for TikZ nodes and big pictures
\newcommand{\bigemoji}[2][1.5cm]{\idiot@art{emoji}{#2}{height=#1}}
% \mascot[<width>]{mood}: panic, happy, scared, sleepy, sick, confused, smug, hungry, thinking
\newcommand{\mascot}[2][2.6cm]{\idiot@art{mascot}{#2}{width=#1}}
\makeatother
% the older figure files place the mascot with \pic{noggin=<mood>}
\tikzset{pics/noggin/.style={code={\node[inner sep=0, transform shape] {\mascot[2.7cm]{#1}};}}}

% ---------------------------------------------------------------------
% Beamer look
% ---------------------------------------------------------------------
\setbeamertemplate{navigation symbols}{}
\setbeamercolor{normal text}{fg=plum, bg=cream}
\setbeamercolor{structure}{fg=plum}
\setbeamercolor{alerted text}{fg=pinkdark}
\setbeamerfont{alerted text}{series=\bfseries}
\setbeamercolor{frametitle}{fg=plum}
\setbeamersize{text margin left=9mm, text margin right=9mm}
\setbeamertemplate{frame begin}{\raggedright} % justified Armenian leaves big gaps

% cream paper with the banner's faint dot grid
\setbeamertemplate{background}{%
  \begin{tikzpicture}
    \fill[cream] (0,0) rectangle (\paperwidth,\paperheight);
    \fill[pattern={Dots[distance=4.4mm, radius=0.28mm]}, pattern color=plum!9!cream]
      (0,0) rectangle (\paperwidth,\paperheight);
  \end{tikzpicture}}

% topic pill, top right of every frame: set by \topic (emoji + short name)
\def\idiottopicemoji{}
\def\idiottopicname{}
\newcommand{\idiotpill}[3][cardwhite]{% fill, emoji, text
  \tikz[baseline=(p.base)]{\node[fill=#1, draw=plum, line width=0.7pt, rounded corners=2.6mm,
      inner xsep=2.2mm, inner ysep=1.1mm, font=\mediumfont\footnotesize, text=plum] (p)
      {\ifx\relax#2\relax\else\emoji[1.05em]{#2}\hspace{1.2mm}\fi#3};
    \begin{scope}[on background layer]
      \fill[plum, rounded corners=2.6mm] ([xshift=0.7mm, yshift=-0.7mm]p.south west) rectangle ([xshift=0.7mm, yshift=-0.7mm]p.north east);
    \end{scope}}}
% \chip{emoji}{text}: inline sticker label, same as the banner's topic chips
\newcommand{\chip}[3][cardwhite]{\idiotpill[#1]{#2}{#3}}

\setbeamertemplate{frametitle}{%
  \vskip 3.5mm
  \begin{minipage}[t]{0.66\paperwidth}
    \raggedright{\blackfont\fontsize{17.5}{20.5}\selectfont\insertframetitle}\par
    \vskip 1.3mm
    \tikz\fill[accent, rounded corners=0.6mm] (0,0) rectangle (1.5cm,1.1mm);
  \end{minipage}%
  \hfill
  \ifx\idiottopicname\empty\else
    \raisebox{-1mm}{\idiotpill{\idiottopicemoji}{\idiottopicname}}%
  \fi
  \vskip 1mm}

\setbeamertemplate{footline}{%
  \hbox to \paperwidth{\hskip 9mm
    {\mediumfont\fontsize{7}{8}\selectfont\color{plum!55!cream}Ապուշ ուղեղը · \insertshorttitle}\hfill
    {\blackfont\fontsize{8}{9}\selectfont\color{plum!70!cream}\insertframenumber}\hskip 9mm}
  \vskip 3mm}

\setbeamertemplate{itemize item}{\tikz[baseline=-0.65ex]\draw[plum, fill=accent, line width=0.6pt] (0,0) circle (0.95mm);}
\setbeamertemplate{itemize subitem}{\tikz[baseline=-0.65ex]\fill[plum!70] (0,0) circle (0.6mm);}
\setbeamertemplate{enumerate item}{\tikz[baseline=(n.base)]\node[fill=accent, draw=plum, line width=0.6pt,
  circle, inner sep=0.5mm, minimum size=4.2mm, font=\blackfont\scriptsize] (n) {\insertenumlabel};}
\setlength{\leftmargini}{5mm}
\setbeamertemplate{itemize/enumerate body begin}{\setlength{\itemsep}{1.6mm}}
\setbeamertemplate{itemize/enumerate subbody begin}{\setlength{\itemsep}{0.6mm}\small}

% ---------------------------------------------------------------------
% Structure: parts (full-colour openers) and topics (book sections)
% ---------------------------------------------------------------------
% \topic{<emoji>}{<short name>}: starts a book section; the pill shows on every following frame
\newcommand{\topic}[2]{\subsection{#2}\gdef\idiottopicemoji{#1}\gdef\idiottopicname{#2}}
\newcommand{\notopic}{\gdef\idiottopicemoji{}\gdef\idiottopicname{}}

% \partopener{<label>}{<title>}{<subtitle>}{<mood>}{<colour>}
\newcommand{\partopener}[5]{%
  \section{#2}\notopic\partaccent{#5}%
  \begin{frame}[plain]
    \begin{tikzpicture}[remember picture, overlay]
      \fill[#5] (current page.south west) rectangle (current page.north east);
      \fill[pattern={Dots[distance=4.4mm, radius=0.28mm]}, pattern color=plum!12!#5]
        (current page.south west) rectangle (current page.north east);
      \fill[cardwhite, opacity=0.45] ([xshift=-3.3cm, yshift=-0.2cm]current page.east) circle (3.1cm);
      \node[anchor=north west, inner sep=0] at ([xshift=10mm, yshift=-12mm]current page.north west)
        {\idiotpill{}{#1}};
      \node[anchor=west, text=plum, align=left, text width=9.2cm, inner sep=0]
        (t) at ([xshift=10mm, yshift=5mm]current page.west)
        {\titlefont\fontsize{38}{40}\selectfont #2};
      \node[anchor=north west, text=plum, text width=9.2cm, inner sep=0, align=left]
        at ([yshift=-5mm]t.south west) {\mediumfont\large #3};
      \node at ([xshift=-3.3cm, yshift=-0.2cm]current page.east) {\mascot[4.6cm]{#4}};
    \end{tikzpicture}
  \end{frame}}

% \topicopener{<emoji>}{<short name>}{<book's section title>}{<what it is about>}{<mood>}
\newcommand{\topicopener}[5]{%
  \topic{#1}{#2}%
  \begin{frame}[plain]
    \begin{tikzpicture}[remember picture, overlay]
      \fill[accent!55!cream] (current page.south west) rectangle (current page.north east);
      \fill[pattern={Dots[distance=4.4mm, radius=0.28mm]}, pattern color=plum!10!accent!55!cream]
        (current page.south west) rectangle (current page.north east);
      \node[anchor=north west, inner sep=0] at ([xshift=10mm, yshift=-12mm]current page.north west)
        {\idiotpill{#1}{#2}};
      \node[anchor=west, text=plum, text width=9.6cm, align=left, inner sep=0]
        (t) at ([xshift=10mm, yshift=6mm]current page.west)
        {\blackfont\fontsize{24}{28}\selectfont «#3»};
      \node[anchor=north west, text=plum, text width=9.6cm, align=left, inner sep=0]
        at ([yshift=-4mm]t.south west) {\mediumfont\large #4};
      \node at ([xshift=-3.0cm, yshift=-0.4cm]current page.east) {\mascot[4.2cm]{#5}};
      \node[rotate=-12] at ([xshift=-1.4cm, yshift=2.3cm]current page.east) {\bigemoji[1.5cm]{#1}};
    \end{tikzpicture}
  \end{frame}}

% ---------------------------------------------------------------------
% Boxes (sticker style: plum outline + hard plum shadow, like the banner)
% ---------------------------------------------------------------------
\tcbset{
  sticker/.style={enhanced, arc=3mm, boxrule=0.9pt, colframe=plum, colback=cardwhite, halign=left,
    left=3mm, right=3mm, top=2mm, bottom=2mm, fontupper=\normalsize, coltext=plum,
    shadow={1.1mm}{-1.1mm}{0mm}{fill=plum}},
  tagged/.style={sticker, colback=#1!45!cream, colbacktitle=#1, coltitle=plum,
    fonttitle=\blackfont\footnotesize, top=3mm,
    attach boxed title to top left={xshift=4mm, yshift*=-2.4mm},
    boxed title style={arc=2.2mm, boxrule=0.9pt, colframe=plum, left=1.5mm, right=1.5mm, top=0.6mm, bottom=0.6mm}},
}
\newtcolorbox{card}[1][]{sticker, #1}
\newtcolorbox{tryit}[1][]{tagged=mint, title={\emoji[1em]{testtube}\ Փորձիր հիմա}, #1}
\newtcolorbox{poll}[1][]{tagged=butter, title={\emoji[1em]{hand}\ Հարցում}, #1}
\newtcolorbox{chatbox}[1][]{tagged=lavender, title={\emoji[1em]{speech}\ Գրիր չաթում}, #1}
\newtcolorbox{scienceupdate}[1][]{tagged=peach, fontupper=\small, title={\emoji[1em]{microscope}\ Գիտությունն այսօր}, #1}
\newtcolorbox{bookquote}[1][]{sticker, colback=cardwhite, fontupper=\mediumfont, left=4mm,
  borderline west={1.4mm}{0pt}{accent},
  before upper={{\blackfont\footnotesize\emoji[1em]{book}\ Բըրնեթը գրում է}\par\smallskip}, #1}

% Brain.exe-style popup window: \begin{popup}{<window title>} ... \end{popup}
\newcommand{\winbuttons}{%
  \tikz[baseline=-0.6ex, line width=0.6pt]{
    \foreach \i/\f in {0/cardwhite, 1/cardwhite, 2/blush}{
      \filldraw[fill=\f, draw=plum, rounded corners=0.6mm] (\i*4.3mm,-1.6mm) rectangle ++(3.3mm,3.3mm);}
    \draw[plum] (0.8mm,-0.8mm) -- (2.5mm,-0.8mm);
    \draw[plum] (5.1mm,-0.8mm) rectangle (6.8mm,0.9mm);
    \draw[plum] (9.4mm,-0.9mm) -- (11.1mm,0.8mm) (9.4mm,0.8mm) -- (11.1mm,-0.9mm);}}
\newtcolorbox{popup}[2][]{sticker, colback=cardwhite, colbacktitle=brainpink, coltitle=plum,
  fonttitle=\blackfont\small, title={\emoji[1.1em]{warning}\hspace{1.5mm}#2\hfill\winbuttons},
  titlerule=0.9pt, toptitle=0.8mm, bottomtitle=0.8mm, #1}

% Myth or fact: \mythorfact{<statement>}{<verdict>}{<colour>}{<explanation>}; verdict shows on slide 2
\newcommand{\mythorfact}[4]{%
  \begin{center}
  \begin{tikzpicture}
    \node[draw=plum, line width=0.9pt, fill=cardwhite, rounded corners=4mm, inner xsep=8mm, inner ysep=6mm,
          text width=10cm, align=center, font=\blackfont\Large, text=plum] (mf) {«#1»};
    \begin{scope}[on background layer]
      \fill[plum, rounded corners=4mm] ([xshift=1.2mm, yshift=-1.2mm]mf.south west) rectangle ([xshift=1.2mm, yshift=-1.2mm]mf.north east);
    \end{scope}
    \node[anchor=west] at ([xshift=6mm]mf.north west) {\idiotpill[butter]{question}{Առասպել, թե՞ փաստ}};
    \onslide<2->{\node[draw=#3, text=#3, line width=2pt, rounded corners=2mm, rotate=-10, font=\blackfont\LARGE,
      inner sep=4pt, fill=cardwhite] at ([xshift=-8mm, yshift=1mm]mf.south east) {#2};}
  \end{tikzpicture}
  \end{center}
  \onslide<2->{#4}}

% Big question with the confused mascot: \bigquestion{<question>}{<small line>}
\newcommand{\bigquestion}[2]{%
  \vfill
  \begin{center}
  \begin{tikzpicture}
    \node[text width=9.6cm, align=center, font=\blackfont\fontsize{24}{29}\selectfont, text=plum] (q) {#1};
    \node[below=4mm of q, text width=9.6cm, align=center, font=\mediumfont\large, text=softgray] {#2};
    \node[anchor=east] at ([xshift=-2mm]q.west) {\mascot[2.6cm]{confused}};
  \end{tikzpicture}
  \end{center}
  \vfill}

% Small source line at the bottom of a frame: \source{Loftus & Palmer, 1974}
\newcommand{\source}[1]{%
  \begin{tikzpicture}[remember picture, overlay]
    \node[anchor=south east, font=\fontsize{6.5}{7.5}\selectfont, text=plum!60!cream, inner sep=0]
      at ([xshift=-17mm, yshift=3.4mm]current page.south east) {Աղբյուր՝ #1};
  \end{tikzpicture}}

% Speech bubbles (handwritten) and the overlay helper for TikZ:
% \node[visible on=<2->] ... keeps the layout but hides the node until slide 2.
\tikzset{
  say/.style={rectangle callout, draw=plum, line width=0.8pt, fill=cardwhite, rounded corners=3pt,
    inner sep=4pt, align=center, font=\handfont\normalsize, text=plum},
  think/.style={cloud callout, cloud puffs=12, draw=plum, line width=0.8pt, fill=cardwhite,
    inner sep=2pt, align=center, font=\handfont\normalsize, text=plum},
  % boxes for diagrams drawn directly in a frame: a style with #1 cannot be defined inside a frame body
  flow/.style={draw=plum, line width=0.8pt, fill=cardwhite, rounded corners=2mm, align=center,
    font=\small, text=plum, inner sep=4pt, text width=#1},
  flow/.default=3cm,
  chatmsg/.style={draw=plum, line width=0.8pt, rounded corners=3mm, inner sep=6pt, text width=#1,
    align=left, font=\small, text=plum},
  chatmsg/.default=7cm,
  invisible/.style={opacity=0, text opacity=0},
  visible on/.style={alt={#1{}{invisible}}},
  alt/.code args={<#1>#2#3}{\alt<#1>{\pgfkeysalso{#2}}{\pgfkeysalso{#3}}},
}

%  Build from the deck folder:  latexmk -lualatex -outdir=build webinarNN.tex
%  Art (emoji, mascot moods, title banner) comes from common/build_art.py.
%
%  All styling lives here, so decks only hold content. Palette and fonts
%  are the banner's (banner/build_banner.py): plum text on cream, pastel
%  accents, Mardoto for text, Glazz for big titles, Twemoji icons and
%  the pink brain mascot (DECISIONS.md #4-#6).
% =====================================================================

% ---------------------------------------------------------------------
% Fonts (LuaLaTeX + fontspec), shared with the banner
% ---------------------------------------------------------------------
\usepackage{fontspec}
\tracinglostchars=3 % a glyph missing from a font stops the build instead of printing nothing
\defaultfontfeatures{Ligatures=TeX}
\def\idiotfontdir{../../banner/assets/fonts/}
\setsansfont{Mardoto}[Path=\idiotfontdir, Extension=.ttf, UprightFont=*-Regular, BoldFont=*-Bold]
\setmainfont{Mardoto}[Path=\idiotfontdir, Extension=.ttf, UprightFont=*-Regular, BoldFont=*-Bold]
\newfontfamily\mediumfont{Mardoto-Medium}[Path=\idiotfontdir, Extension=.ttf]
\newfontfamily\blackfont{Mardoto-Black}[Path=\idiotfontdir, Extension=.ttf]
\newfontfamily\titlefont{GlazzArm-Regular}[Path=\idiotfontdir, Extension=.ttf] % Armenian letters only
\newfontfamily\handfont{Ink Free} % Windows font with Armenian; speech bubbles only
% Mardoto has no italic: emphasis is medium weight in the accent colour instead.
\renewcommand{\emph}[1]{{\mediumfont\color{plum}#1}}
% No hyphenation: split words are hard to read on a shared screen.
\hyphenpenalty=10000
\exhyphenpenalty=10000

% ---------------------------------------------------------------------
% Packages
% ---------------------------------------------------------------------
\usepackage{graphicx}
\usepackage{booktabs}
\usepackage[most]{tcolorbox}
\usepackage{tikz}
\usetikzlibrary{arrows.meta, positioning, calc, shapes.geometric, shapes.callouts, shapes.misc,
  decorations.pathmorphing, decorations.pathreplacing, backgrounds, fit, patterns.meta}
\usepackage{pgfplots}
\pgfplotsset{compat=1.18}
\usepgfplotslibrary{fillbetween}
\graphicspath{{../common/art/}{images/}}

% ---------------------------------------------------------------------
% Colours: the banner palette. Charts with 3+ series use the Armenian
% flag colours (amRed / amBlue / amOrange).
% ---------------------------------------------------------------------
\definecolor{plum}{HTML}{4A2340}
\definecolor{brainpink}{HTML}{F7A8C4}
\definecolor{fold}{HTML}{E57FA6}
\definecolor{cream}{HTML}{FFF6EC}
\definecolor{cardwhite}{HTML}{FFFDF9}
\definecolor{peach}{HTML}{FFD9C2}
\definecolor{lavender}{HTML}{D9C8F5}
\definecolor{mint}{HTML}{B5E8D5}
\definecolor{butter}{HTML}{FFE8A3}
\definecolor{blush}{HTML}{FF8FA8}
\definecolor{amRed}{HTML}{D90012}
\definecolor{amBlue}{HTML}{0033A0}
\definecolor{amOrange}{HTML}{F2A800}
% darker partners for text on cream
\colorlet{mintdark}{mint!45!black!80!green}
\definecolor{mintdark}{HTML}{2E7D62}
\definecolor{lavenderdark}{HTML}{5B3FA0}
\definecolor{peachdark}{HTML}{B5532A}
\definecolor{pinkdark}{HTML}{C2185B}
% names used by the figure files
\colorlet{ink}{plum}
\colorlet{paper}{cream}
\colorlet{mist}{lavender!40!cream}
\definecolor{night}{HTML}{2B2150}
\colorlet{softgray}{plum!60!cream}
% accent of the current part; \partaccent{<colour>} switches it
\colorlet{accent}{butter}
\newcommand{\partaccent}[1]{\colorlet{accent}{#1}}

% ---------------------------------------------------------------------
% Art: \emoji, \mascot (pages of common/art/*.pdf, see build_art.py)
% ---------------------------------------------------------------------
\makeatletter
\newcommand{\artpage}[3]{\expandafter\def\csname idiot@art@#1@#2\endcsname{#3}}
% generated by build_art.py - do not edit
\artpage{emoji}{brain}{1}
\artpage{emoji}{warning}{2}
\artpage{emoji}{rollercoaster}{3}
\artpage{emoji}{doughnut}{4}
\artpage{emoji}{sleeping}{5}
\artpage{emoji}{scream}{6}
\artpage{emoji}{thought}{7}
\artpage{emoji}{question}{8}
\artpage{emoji}{cyclone}{9}
\artpage{emoji}{calendar}{10}
\artpage{emoji}{clock630}{11}
\artpage{emoji}{computer}{12}
\artpage{emoji}{receipt}{13}
\artpage{emoji}{gift}{14}
\artpage{emoji}{lizard}{15}
\artpage{emoji}{boss}{16}
\artpage{emoji}{keyboard}{17}
\artpage{emoji}{wastebasket}{18}
\artpage{emoji}{batterylow}{19}
\artpage{emoji}{skull}{20}
\artpage{emoji}{car}{21}
\artpage{emoji}{ship}{22}
\artpage{emoji}{airplane}{23}
\artpage{emoji}{nauseated}{24}
\artpage{emoji}{vomit}{25}
\artpage{emoji}{eye}{26}
\artpage{emoji}{ear}{27}
\artpage{emoji}{tree}{28}
\artpage{emoji}{wave}{29}
\artpage{emoji}{book}{30}
\artpage{emoji}{walking}{31}
\artpage{emoji}{dizzy}{32}
\artpage{emoji}{cake}{33}
\artpage{emoji}{icecream}{34}
\artpage{emoji}{potato}{35}
\artpage{emoji}{apple}{36}
\artpage{emoji}{burger}{37}
\artpage{emoji}{salad}{38}
\artpage{emoji}{clock12}{39}
\artpage{emoji}{forkknife}{40}
\artpage{emoji}{pasta}{41}
\artpage{emoji}{fire}{42}
\artpage{emoji}{drool}{43}
\artpage{emoji}{pizza}{44}
\artpage{emoji}{bed}{45}
\artpage{emoji}{personbed}{46}
\artpage{emoji}{moon}{47}
\artpage{emoji}{sun}{48}
\artpage{emoji}{zzz}{49}
\artpage{emoji}{dolphin}{50}
\artpage{emoji}{jellyfish}{51}
\artpage{emoji}{rat}{52}
\artpage{emoji}{chipmunk}{53}
\artpage{emoji}{alarm}{54}
\artpage{emoji}{phone}{55}
\artpage{emoji}{ghost}{56}
\artpage{emoji}{alien}{57}
\artpage{emoji}{yawn}{58}
\artpage{emoji}{baby}{59}
\artpage{emoji}{construction}{60}
\artpage{emoji}{broom}{61}
\artpage{emoji}{thermometer}{62}
\artpage{emoji}{cold}{63}
\artpage{emoji}{monkey}{64}
\artpage{emoji}{tiger}{65}
\artpage{emoji}{axe}{66}
\artpage{emoji}{balloon}{67}
\artpage{emoji}{heart}{68}
\artpage{emoji}{anatomicalheart}{69}
\artpage{emoji}{lungs}{70}
\artpage{emoji}{sweat}{71}
\artpage{emoji}{eyes}{72}
\artpage{emoji}{butterfly}{73}
\artpage{emoji}{runner}{74}
\artpage{emoji}{boxing}{75}
\artpage{emoji}{station}{76}
\artpage{emoji}{taxi}{77}
\artpage{emoji}{bell}{78}
\artpage{emoji}{couch}{79}
\artpage{emoji}{muscle}{80}
\artpage{emoji}{toilet}{81}
\artpage{emoji}{floppy}{82}
\artpage{emoji}{door}{83}
\artpage{emoji}{coffee}{84}
\artpage{emoji}{diving}{85}
\artpage{emoji}{wine}{86}
\artpage{emoji}{beer}{87}
\artpage{emoji}{dove}{88}
\artpage{emoji}{puzzle}{89}
\artpage{emoji}{fish}{90}
\artpage{emoji}{trophy}{91}
\artpage{emoji}{detective}{92}
\artpage{emoji}{judge}{93}
\artpage{emoji}{cheese}{94}
\artpage{emoji}{film}{95}
\artpage{emoji}{magnifier}{96}
\artpage{emoji}{key}{97}
\artpage{emoji}{firework}{98}
\artpage{emoji}{house}{99}
\artpage{emoji}{gloves}{100}
\artpage{emoji}{trafficlight}{101}
\artpage{emoji}{stop}{102}
\artpage{emoji}{mute}{103}
\artpage{emoji}{videogame}{104}
\artpage{emoji}{mirror}{105}
\artpage{emoji}{star}{106}
\artpage{emoji}{hospital}{107}
\artpage{emoji}{pill}{108}
\artpage{emoji}{camera}{109}
\artpage{emoji}{tv}{110}
\artpage{emoji}{crown}{111}
\artpage{emoji}{nerd}{112}
\artpage{emoji}{lock}{113}
\artpage{emoji}{musical}{114}
\artpage{emoji}{chips}{115}
\artpage{emoji}{check}{116}
\artpage{emoji}{cross}{117}
\artpage{emoji}{hand}{118}
\artpage{emoji}{wavehand}{119}
\artpage{emoji}{speech}{120}
\artpage{emoji}{pencil}{121}
\artpage{emoji}{memo}{122}
\artpage{emoji}{monocle}{123}
\artpage{emoji}{thinking}{124}
\artpage{emoji}{grin}{125}
\artpage{emoji}{sweatsmile}{126}
\artpage{emoji}{facepalm}{127}
\artpage{emoji}{popcorn}{128}
\artpage{emoji}{bulb}{129}
\artpage{emoji}{sparkles}{130}
\artpage{emoji}{party}{131}
\artpage{emoji}{chart}{132}
\artpage{emoji}{testtube}{133}
\artpage{emoji}{people}{134}
\artpage{emoji}{pin}{135}
\artpage{emoji}{hourglass}{136}
\artpage{emoji}{microscope}{137}
\artpage{emoji}{books}{138}
\artpage{emoji}{megaphone}{139}
\artpage{emoji}{student}{140}
\artpage{emoji}{zany}{141}
\artpage{emoji}{hundred}{142}
\artpage{emoji}{link}{143}
\artpage{emoji}{target}{144}
\artpage{emoji}{clover}{145}
\artpage{emoji}{ufo}{146}
\artpage{emoji}{spider}{147}
\artpage{emoji}{spiderweb}{148}
\artpage{emoji}{snake}{149}
\artpage{emoji}{microphone}{150}
\artpage{emoji}{clown}{151}
\artpage{emoji}{pumpkin}{152}
\artpage{emoji}{knife}{153}
\artpage{emoji}{bat}{154}
\artpage{emoji}{blackcat}{155}
\artpage{emoji}{crossedfingers}{156}
\artpage{emoji}{ladder}{157}
\artpage{emoji}{thumbsup}{158}
\artpage{emoji}{thumbsdown}{159}
\artpage{emoji}{rocket}{160}
\artpage{emoji}{satellite}{161}
\artpage{emoji}{newspaper}{162}
\artpage{emoji}{chocolate}{163}
\artpage{emoji}{syringe}{164}
\artpage{emoji}{dice}{165}
\artpage{emoji}{hauntedhouse}{166}
\artpage{emoji}{flushed}{167}
\artpage{emoji}{fearful}{168}
\artpage{emoji}{graduation}{169}
\artpage{emoji}{teacher}{170}
\artpage{emoji}{scientist}{171}
\artpage{emoji}{chess}{172}
\artpage{emoji}{abacus}{173}
\artpage{emoji}{ruler}{174}
\artpage{emoji}{jeans}{175}
\artpage{emoji}{glasses}{176}
\artpage{emoji}{dna}{177}
\artpage{emoji}{loudspeaker}{178}
\artpage{emoji}{speakinghead}{179}
\artpage{emoji}{seedling}{180}
\artpage{emoji}{weightlifter}{181}
\artpage{emoji}{giraffe}{182}
\artpage{emoji}{mouse}{183}
\artpage{emoji}{exploding}{184}
\artpage{emoji}{owl}{185}
\artpage{emoji}{clipboard}{186}
\artpage{emoji}{numbers}{187}
\artpage{emoji}{medal}{188}
\artpage{emoji}{rose}{189}
\artpage{emoji}{nose}{190}
\artpage{emoji}{tongue}{191}
\artpage{emoji}{yum}{192}
\artpage{emoji}{bread}{193}
\artpage{emoji}{lemon}{194}
\artpage{emoji}{chili}{195}
\artpage{emoji}{salt}{196}
\artpage{emoji}{candy}{197}
\artpage{emoji}{speaker}{198}
\artpage{emoji}{headphones}{199}
\artpage{emoji}{feather}{200}
\artpage{emoji}{gorilla}{201}
\artpage{emoji}{basketball}{202}
\artpage{emoji}{magicwand}{203}
\artpage{emoji}{tophat}{204}
\artpage{emoji}{cloud}{205}
\artpage{emoji}{plug}{206}
\artpage{emoji}{clinking}{207}
\artpage{emoji}{dog}{208}
\artpage{emoji}{rainbow}{209}
\artpage{emoji}{slightsmile}{210}
\artpage{emoji}{palette}{211}
\artpage{emoji}{masks}{212}
\artpage{emoji}{angry}{213}
\artpage{emoji}{rage}{214}
\artpage{emoji}{steam}{215}
\artpage{emoji}{bomb}{216}
\artpage{emoji}{volcano}{217}
\artpage{emoji}{angersymbol}{218}
\artpage{emoji}{crystalball}{219}
\artpage{emoji}{railway}{220}
\artpage{emoji}{worker}{221}
\artpage{emoji}{laughing}{222}
\artpage{emoji}{rofl}{223}
\artpage{emoji}{banana}{224}
\artpage{emoji}{mountain}{225}
\artpage{emoji}{moneybag}{226}
\artpage{emoji}{lightning}{227}
\artpage{emoji}{joker}{228}
\artpage{emoji}{finishflag}{229}
\artpage{emoji}{hug}{230}
\artpage{emoji}{handshake}{231}
\artpage{emoji}{lying}{232}
\artpage{emoji}{neutral}{233}
\artpage{emoji}{carrot}{234}
\artpage{emoji}{wood}{235}
\artpage{emoji}{brokenheart}{236}
\artpage{emoji}{loveletter}{237}
\artpage{emoji}{couple}{238}
\artpage{emoji}{kiss}{239}
\artpage{emoji}{cowboy}{240}
\artpage{emoji}{guitar}{241}
\artpage{emoji}{family}{242}
\artpage{emoji}{police}{243}
\artpage{emoji}{sheep}{244}
\artpage{emoji}{clapping}{245}
\artpage{emoji}{imp}{246}
\artpage{emoji}{ballot}{247}
\artpage{emoji}{smileheart}{248}
\artpage{emoji}{dollar}{249}
\artpage{emoji}{earth}{250}
\artpage{emoji}{raincloud}{251}
\artpage{emoji}{sunbehindcloud}{252}
\artpage{emoji}{pensive}{253}
\artpage{emoji}{cry}{254}
\artpage{emoji}{bandage}{255}
\artpage{emoji}{stethoscope}{256}
\artpage{emoji}{cigarette}{257}
\artpage{emoji}{slotmachine}{258}
\artpage{emoji}{siren}{259}
\artpage{emoji}{sunrise}{260}
\artpage{emoji}{healthworker}{261}
\artpage{emoji}{hearnoevil}{262}
\artpage{emoji}{anchor}{263}
\artpage{emoji}{candle}{264}
\artpage{emoji}{telephone}{265}
\artpage{emoji}{battery}{266}
\artpage{emoji}{weary}{267}
\artpage{emoji}{lotus}{268}
\artpage{mascot}{panic}{1}
\artpage{mascot}{happy}{2}
\artpage{mascot}{scared}{3}
\artpage{mascot}{sleepy}{4}
\artpage{mascot}{sick}{5}
\artpage{mascot}{confused}{6}
\artpage{mascot}{smug}{7}
\artpage{mascot}{hungry}{8}
\artpage{mascot}{thinking}{9}
\artpage{mascot}{sad}{10}
\artpage{mascot}{angry}{11}

\newcommand{\idiot@art}[3]{% kind, name, size options
  \ifcsname idiot@art@#1@#2\endcsname
    \includegraphics[#3, page=\csname idiot@art@#1@#2\endcsname]{#1.pdf}%
  \else
    \PackageError{idiot}{Unknown #1 '#2'; add it to build_art.py and rebuild the art}{}%
  \fi}
% \emoji[<height>]{name}: inline icon, sits on the text baseline like a letter
\newcommand{\emoji}[2][1.15em]{\raisebox{-0.2em}{\idiot@art{emoji}{#2}{height=#1}}}
% \bigemoji[<height>]{name}: no baseline shift, for TikZ nodes and big pictures
\newcommand{\bigemoji}[2][1.5cm]{\idiot@art{emoji}{#2}{height=#1}}
% \mascot[<width>]{mood}: panic, happy, scared, sleepy, sick, confused, smug, hungry, thinking
\newcommand{\mascot}[2][2.6cm]{\idiot@art{mascot}{#2}{width=#1}}
\makeatother
% the older figure files place the mascot with \pic{noggin=<mood>}
\tikzset{pics/noggin/.style={code={\node[inner sep=0, transform shape] {\mascot[2.7cm]{#1}};}}}

% ---------------------------------------------------------------------
% Beamer look
% ---------------------------------------------------------------------
\setbeamertemplate{navigation symbols}{}
\setbeamercolor{normal text}{fg=plum, bg=cream}
\setbeamercolor{structure}{fg=plum}
\setbeamercolor{alerted text}{fg=pinkdark}
\setbeamerfont{alerted text}{series=\bfseries}
\setbeamercolor{frametitle}{fg=plum}
\setbeamersize{text margin left=9mm, text margin right=9mm}
\setbeamertemplate{frame begin}{\raggedright} % justified Armenian leaves big gaps

% cream paper with the banner's faint dot grid
\setbeamertemplate{background}{%
  \begin{tikzpicture}
    \fill[cream] (0,0) rectangle (\paperwidth,\paperheight);
    \fill[pattern={Dots[distance=4.4mm, radius=0.28mm]}, pattern color=plum!9!cream]
      (0,0) rectangle (\paperwidth,\paperheight);
  \end{tikzpicture}}

% topic pill, top right of every frame: set by \topic (emoji + short name)
\def\idiottopicemoji{}
\def\idiottopicname{}
\newcommand{\idiotpill}[3][cardwhite]{% fill, emoji, text
  \tikz[baseline=(p.base)]{\node[fill=#1, draw=plum, line width=0.7pt, rounded corners=2.6mm,
      inner xsep=2.2mm, inner ysep=1.1mm, font=\mediumfont\footnotesize, text=plum] (p)
      {\ifx\relax#2\relax\else\emoji[1.05em]{#2}\hspace{1.2mm}\fi#3};
    \begin{scope}[on background layer]
      \fill[plum, rounded corners=2.6mm] ([xshift=0.7mm, yshift=-0.7mm]p.south west) rectangle ([xshift=0.7mm, yshift=-0.7mm]p.north east);
    \end{scope}}}
% \chip{emoji}{text}: inline sticker label, same as the banner's topic chips
\newcommand{\chip}[3][cardwhite]{\idiotpill[#1]{#2}{#3}}

\setbeamertemplate{frametitle}{%
  \vskip 3.5mm
  \begin{minipage}[t]{0.66\paperwidth}
    \raggedright{\blackfont\fontsize{17.5}{20.5}\selectfont\insertframetitle}\par
    \vskip 1.3mm
    \tikz\fill[accent, rounded corners=0.6mm] (0,0) rectangle (1.5cm,1.1mm);
  \end{minipage}%
  \hfill
  \ifx\idiottopicname\empty\else
    \raisebox{-1mm}{\idiotpill{\idiottopicemoji}{\idiottopicname}}%
  \fi
  \vskip 1mm}

\setbeamertemplate{footline}{%
  \hbox to \paperwidth{\hskip 9mm
    {\mediumfont\fontsize{7}{8}\selectfont\color{plum!55!cream}Ապուշ ուղեղը · \insertshorttitle}\hfill
    {\blackfont\fontsize{8}{9}\selectfont\color{plum!70!cream}\insertframenumber}\hskip 9mm}
  \vskip 3mm}

\setbeamertemplate{itemize item}{\tikz[baseline=-0.65ex]\draw[plum, fill=accent, line width=0.6pt] (0,0) circle (0.95mm);}
\setbeamertemplate{itemize subitem}{\tikz[baseline=-0.65ex]\fill[plum!70] (0,0) circle (0.6mm);}
\setbeamertemplate{enumerate item}{\tikz[baseline=(n.base)]\node[fill=accent, draw=plum, line width=0.6pt,
  circle, inner sep=0.5mm, minimum size=4.2mm, font=\blackfont\scriptsize] (n) {\insertenumlabel};}
\setlength{\leftmargini}{5mm}
\setbeamertemplate{itemize/enumerate body begin}{\setlength{\itemsep}{1.6mm}}
\setbeamertemplate{itemize/enumerate subbody begin}{\setlength{\itemsep}{0.6mm}\small}

% ---------------------------------------------------------------------
% Structure: parts (full-colour openers) and topics (book sections)
% ---------------------------------------------------------------------
% \topic{<emoji>}{<short name>}: starts a book section; the pill shows on every following frame
\newcommand{\topic}[2]{\subsection{#2}\gdef\idiottopicemoji{#1}\gdef\idiottopicname{#2}}
\newcommand{\notopic}{\gdef\idiottopicemoji{}\gdef\idiottopicname{}}

% \partopener{<label>}{<title>}{<subtitle>}{<mood>}{<colour>}
\newcommand{\partopener}[5]{%
  \section{#2}\notopic\partaccent{#5}%
  \begin{frame}[plain]
    \begin{tikzpicture}[remember picture, overlay]
      \fill[#5] (current page.south west) rectangle (current page.north east);
      \fill[pattern={Dots[distance=4.4mm, radius=0.28mm]}, pattern color=plum!12!#5]
        (current page.south west) rectangle (current page.north east);
      \fill[cardwhite, opacity=0.45] ([xshift=-3.3cm, yshift=-0.2cm]current page.east) circle (3.1cm);
      \node[anchor=north west, inner sep=0] at ([xshift=10mm, yshift=-12mm]current page.north west)
        {\idiotpill{}{#1}};
      \node[anchor=west, text=plum, align=left, text width=9.2cm, inner sep=0]
        (t) at ([xshift=10mm, yshift=5mm]current page.west)
        {\titlefont\fontsize{38}{40}\selectfont #2};
      \node[anchor=north west, text=plum, text width=9.2cm, inner sep=0, align=left]
        at ([yshift=-5mm]t.south west) {\mediumfont\large #3};
      \node at ([xshift=-3.3cm, yshift=-0.2cm]current page.east) {\mascot[4.6cm]{#4}};
    \end{tikzpicture}
  \end{frame}}

% \topicopener{<emoji>}{<short name>}{<book's section title>}{<what it is about>}{<mood>}
\newcommand{\topicopener}[5]{%
  \topic{#1}{#2}%
  \begin{frame}[plain]
    \begin{tikzpicture}[remember picture, overlay]
      \fill[accent!55!cream] (current page.south west) rectangle (current page.north east);
      \fill[pattern={Dots[distance=4.4mm, radius=0.28mm]}, pattern color=plum!10!accent!55!cream]
        (current page.south west) rectangle (current page.north east);
      \node[anchor=north west, inner sep=0] at ([xshift=10mm, yshift=-12mm]current page.north west)
        {\idiotpill{#1}{#2}};
      \node[anchor=west, text=plum, text width=9.6cm, align=left, inner sep=0]
        (t) at ([xshift=10mm, yshift=6mm]current page.west)
        {\blackfont\fontsize{24}{28}\selectfont «#3»};
      \node[anchor=north west, text=plum, text width=9.6cm, align=left, inner sep=0]
        at ([yshift=-4mm]t.south west) {\mediumfont\large #4};
      \node at ([xshift=-3.0cm, yshift=-0.4cm]current page.east) {\mascot[4.2cm]{#5}};
      \node[rotate=-12] at ([xshift=-1.4cm, yshift=2.3cm]current page.east) {\bigemoji[1.5cm]{#1}};
    \end{tikzpicture}
  \end{frame}}

% ---------------------------------------------------------------------
% Boxes (sticker style: plum outline + hard plum shadow, like the banner)
% ---------------------------------------------------------------------
\tcbset{
  sticker/.style={enhanced, arc=3mm, boxrule=0.9pt, colframe=plum, colback=cardwhite, halign=left,
    left=3mm, right=3mm, top=2mm, bottom=2mm, fontupper=\normalsize, coltext=plum,
    shadow={1.1mm}{-1.1mm}{0mm}{fill=plum}},
  tagged/.style={sticker, colback=#1!45!cream, colbacktitle=#1, coltitle=plum,
    fonttitle=\blackfont\footnotesize, top=3mm,
    attach boxed title to top left={xshift=4mm, yshift*=-2.4mm},
    boxed title style={arc=2.2mm, boxrule=0.9pt, colframe=plum, left=1.5mm, right=1.5mm, top=0.6mm, bottom=0.6mm}},
}
\newtcolorbox{card}[1][]{sticker, #1}
\newtcolorbox{tryit}[1][]{tagged=mint, title={\emoji[1em]{testtube}\ Փորձիր հիմա}, #1}
\newtcolorbox{poll}[1][]{tagged=butter, title={\emoji[1em]{hand}\ Հարցում}, #1}
\newtcolorbox{chatbox}[1][]{tagged=lavender, title={\emoji[1em]{speech}\ Գրիր չաթում}, #1}
\newtcolorbox{scienceupdate}[1][]{tagged=peach, fontupper=\small, title={\emoji[1em]{microscope}\ Գիտությունն այսօր}, #1}
\newtcolorbox{bookquote}[1][]{sticker, colback=cardwhite, fontupper=\mediumfont, left=4mm,
  borderline west={1.4mm}{0pt}{accent},
  before upper={{\blackfont\footnotesize\emoji[1em]{book}\ Բըրնեթը գրում է}\par\smallskip}, #1}

% Brain.exe-style popup window: \begin{popup}{<window title>} ... \end{popup}
\newcommand{\winbuttons}{%
  \tikz[baseline=-0.6ex, line width=0.6pt]{
    \foreach \i/\f in {0/cardwhite, 1/cardwhite, 2/blush}{
      \filldraw[fill=\f, draw=plum, rounded corners=0.6mm] (\i*4.3mm,-1.6mm) rectangle ++(3.3mm,3.3mm);}
    \draw[plum] (0.8mm,-0.8mm) -- (2.5mm,-0.8mm);
    \draw[plum] (5.1mm,-0.8mm) rectangle (6.8mm,0.9mm);
    \draw[plum] (9.4mm,-0.9mm) -- (11.1mm,0.8mm) (9.4mm,0.8mm) -- (11.1mm,-0.9mm);}}
\newtcolorbox{popup}[2][]{sticker, colback=cardwhite, colbacktitle=brainpink, coltitle=plum,
  fonttitle=\blackfont\small, title={\emoji[1.1em]{warning}\hspace{1.5mm}#2\hfill\winbuttons},
  titlerule=0.9pt, toptitle=0.8mm, bottomtitle=0.8mm, #1}

% Myth or fact: \mythorfact{<statement>}{<verdict>}{<colour>}{<explanation>}; verdict shows on slide 2
\newcommand{\mythorfact}[4]{%
  \begin{center}
  \begin{tikzpicture}
    \node[draw=plum, line width=0.9pt, fill=cardwhite, rounded corners=4mm, inner xsep=8mm, inner ysep=6mm,
          text width=10cm, align=center, font=\blackfont\Large, text=plum] (mf) {«#1»};
    \begin{scope}[on background layer]
      \fill[plum, rounded corners=4mm] ([xshift=1.2mm, yshift=-1.2mm]mf.south west) rectangle ([xshift=1.2mm, yshift=-1.2mm]mf.north east);
    \end{scope}
    \node[anchor=west] at ([xshift=6mm]mf.north west) {\idiotpill[butter]{question}{Առասպել, թե՞ փաստ}};
    \onslide<2->{\node[draw=#3, text=#3, line width=2pt, rounded corners=2mm, rotate=-10, font=\blackfont\LARGE,
      inner sep=4pt, fill=cardwhite] at ([xshift=-8mm, yshift=1mm]mf.south east) {#2};}
  \end{tikzpicture}
  \end{center}
  \onslide<2->{#4}}

% Big question with the confused mascot: \bigquestion{<question>}{<small line>}
\newcommand{\bigquestion}[2]{%
  \vfill
  \begin{center}
  \begin{tikzpicture}
    \node[text width=9.6cm, align=center, font=\blackfont\fontsize{24}{29}\selectfont, text=plum] (q) {#1};
    \node[below=4mm of q, text width=9.6cm, align=center, font=\mediumfont\large, text=softgray] {#2};
    \node[anchor=east] at ([xshift=-2mm]q.west) {\mascot[2.6cm]{confused}};
  \end{tikzpicture}
  \end{center}
  \vfill}

% Small source line at the bottom of a frame: \source{Loftus & Palmer, 1974}
\newcommand{\source}[1]{%
  \begin{tikzpicture}[remember picture, overlay]
    \node[anchor=south east, font=\fontsize{6.5}{7.5}\selectfont, text=plum!60!cream, inner sep=0]
      at ([xshift=-17mm, yshift=3.4mm]current page.south east) {Աղբյուր՝ #1};
  \end{tikzpicture}}

% Speech bubbles (handwritten) and the overlay helper for TikZ:
% \node[visible on=<2->] ... keeps the layout but hides the node until slide 2.
\tikzset{
  say/.style={rectangle callout, draw=plum, line width=0.8pt, fill=cardwhite, rounded corners=3pt,
    inner sep=4pt, align=center, font=\handfont\normalsize, text=plum},
  think/.style={cloud callout, cloud puffs=12, draw=plum, line width=0.8pt, fill=cardwhite,
    inner sep=2pt, align=center, font=\handfont\normalsize, text=plum},
  % boxes for diagrams drawn directly in a frame: a style with #1 cannot be defined inside a frame body
  flow/.style={draw=plum, line width=0.8pt, fill=cardwhite, rounded corners=2mm, align=center,
    font=\small, text=plum, inner sep=4pt, text width=#1},
  flow/.default=3cm,
  chatmsg/.style={draw=plum, line width=0.8pt, rounded corners=3mm, inner sep=6pt, text width=#1,
    align=left, font=\small, text=plum},
  chatmsg/.default=7cm,
  invisible/.style={opacity=0, text opacity=0},
  visible on/.style={alt={#1{}{invisible}}},
  alt/.code args={<#1>#2#3}{\alt<#1>{\pgfkeysalso{#2}}{\pgfkeysalso{#3}}},
}
